\chapter{The trace law and its consequences}
\label{ch17:trace}

The leading short-distance limit of the stress tensor is the traceless
Maxwell tensor. It therefore says nothing by itself about the leading
nonzero term in the trace. That term is smaller by one inverse running
coupling, and its two-point function by two. The perturbative origin of
this phenomenon is the trace anomaly of Collins, Duncan and Joglekar
\cite{CollinsDuncanJoglekar1977}. The all-radius assertion below requires
additional estimates for the selected lattice construction; it does not
follow by substituting a perturbative beta function into a formal identity.

Retain the Ward-normalized tensor $T$ of
Proposition~\ref{ch14:normalized-stress}, the phase $\xi_*$ chosen in
Section~\ref{ch15:selection}, and its relativistic reconstruction.
The latter supplies a positive-cone translation spectrum and tempered
boundary correlations. We prove the stronger quantitative current
transfer needed for the trace asymptotic, including its low-moment
cancellations. The construction and reconstruction in the preceding
chapters do not depend on this refinement.

\section{Centering and the positive spectral measure}
\label{ch17:spectrum}

Put $\tau=\sum_{a=1}^4T_{aa}$, $e_{\rm tr}=S^{\xi_*}(\tau)$ and
\[
             \Theta=\tau-e_{\rm tr}1,\qquad
             G_\Theta(r)=S_2^{\xi_*}(\Theta,\Theta;re_4,0).
\]
The real number $e_{\rm tr}$ is the expectation in the selected phase,
which need not equal the mixed expectation. The unit field is retained
through the common filtration. The fact that inserting the identity leaves correlations unchanged extends to
the boundary correlations, so, on the common invariant word domain,
\[
 \Phi_\Theta(f)=\Phi_\tau(f)
                  -e_{\rm tr}\left(\int f\right)I,
             \qquad \langle\Omega,\Phi_\Theta(f)\Omega\rangle=0.
\]
For completeness, the Euclidean correlation with a unit insertion is
constant in that insertion position. Holomorphic uniqueness on ordered
time regions identifies its continuation with the correlation after
deletion. Boundary values and smearing give the displayed operator
identity on dense word matrix elements. This establishes centering on
the domain, as well as for separated expectations.

Write $p=(\boldsymbol p,E)$, $p^2=E^2-|\boldsymbol p|^2$, and
$\overline V_+=\{E\ge|\boldsymbol p|\}$. Use
$\widehat f_+(p)=\int e^{ip\cdot x}f(x)\,dx$ in the translation
spectral convention, and
$\mathcal W\psi(p)=\int e^{ip\cdot(y-x)}\psi(x,y)\,dx\,dy$.
Then $\mathcal W(\overline f\otimes h)=
\overline{\widehat f_+}\widehat h_+$.
Positivity, translation covariance and temperedness give a unique
positive measure $\rho_\Theta$ of polynomial growth, supported in
$\overline V_+$, with
\begin{align}
 W_2(\Theta,\Theta;\psi)&=\int\mathcal W\psi\,d\rho_\Theta,
 \label{ch17:spectral-pairing}\\
 G_\Theta(r)&=\int e^{-rE}\,d\rho_\Theta(p).
 \label{ch17:laplace}
\end{align}
Here is the direct spectral identification needed for these formulas.
For $v_f=\Phi_\Theta(f)\Omega$, covariance gives
$\widehat h_+(P)v_f=\widehat f_+(P)v_h$ by convolution.
On a compact momentum set choose $\widehat f_+$ nonzero and define
$d\rho_\Theta=|\widehat f_+|^{-2}d\mu_{v_f}$ there. The convolution
identity makes the definitions agree on overlaps. Polarization gives
\eqref{ch17:spectral-pairing} first on product test functions and then on joint
Schwartz test functions. Tempered continuity bounds compact masses polynomially.
The positive cone permits exponential damping, and uniqueness of the
holomorphic boundary value identifies its Laplace transform with the
Euclidean function. The scalar transformation law of $\Theta$ makes
$\rho_\Theta$ Lorentz invariant. In particular $G_\Theta$ is real,
nonnegative and nonincreasing. The expression
$r^8\log^2(1/r)G_\Theta(r)$ need not be monotone.

\section{Comparison at prescribed radii}
\label{ch17:radii}

Let $U_{i,d}:\mathcal M_d^{\rm Max}\to F_d$ be real comparison
isomorphisms, preserving the unit, inclusions and signed coordinate
symmetries. They are evaluated at a radius $r_i\downarrow0$:
$U_{i,d}A$ is inserted at $r_ix$. The actual inverse coupling at a
finer insertion level $k_i$ is denoted by $\beta_i$, and
$b_i=\widehat\beta_{k_i}$ is its reference value. The radius level
$n_i$ precedes $k_i$.

The strengthened comparison, proved below, is as follows. Every prescribed
positive null sequence has a subsequence with these data such that
\begin{equation}
 \beta_i/b_i\to1,\qquad
 \frac{\beta_i}{\log(1/r_i)}\to\kappa_G,
 \label{ch17:running-radius}
\end{equation}
and all fixed separated correlations and their fixed derivatives,
with convergent moving free labels, converge to their Maxwell values.
For pairs without derivatives this comparison is uniform over the
phases, including at nearest lattice-rounded points. Shell integrals are
included. It must use the same field quotient and lattice family as
before. A comparison along one freely chosen scale sequence is weaker.

We explain the radius coverage and the relevant quantitative demands.
Write $M_n=M_n^{\rm cut}=\exp((\log\log\widehat\beta_n)^2)$ and
allow
\[
 r(n,u)=a_nM_n^u,\quad \alpha_0\le u\le\alpha_0+1,\qquad
 k(n)=\min\{m>n:\log(a_n/a_m)\ge(\log\widehat\beta_n)^{48}\}.
\]
At fixed degree bound $D$ and arity bound $N_*$, choose comparison
exponents with
\begin{gather*}
 c<0,\quad\alpha_0-\alpha_1>6,\quad
 \alpha_0-\alpha_2>5-c,\\
 (\alpha_0-\alpha_1)DN_*+c<0,\qquad
 (\alpha_0-\alpha_2)DN_*-\alpha_0-c<0.
\end{gather*}
The common size of the $\alpha_j$ can be increased after their
differences and $c$ have been fixed. The marked comparison required
at every sufficiently late level has, uniformly in $u$, single-insertion
seminorm bounds $CM_n^{(\alpha_0-\alpha_j)D}$ and grid errors tending
to zero. With strip width $M_n^c$ and scaled mesh at most
$M_n^{-\alpha_0}$, the two integration errors are bounded by
\[
 C M_n^{(\alpha_0-\alpha_1)DN_*+c},\qquad
 C M_n^{(\alpha_0-\alpha_2)DN_*-\alpha_0-c}.
\]
They tend to zero. Thus the strip removal and derivative-transport
argument of Section~\ref{ch14:comparison} applies to these windows.
Bounds uniform in the number of copies give phase uniformity by
\eqref{ch15:moment-norm}. The next subsection verifies these window estimates, including their
uniformity across changes of stage, from marked propagation.

The coupling thresholds can be chosen to ensure $L_n\le\log\widehat\beta_n$.
Since $M_n$ increases and $L_n\le M_{n+1}$ eventually,
$r(n,\alpha_0)\le r(n+1,\alpha_0+1)$. The intervals therefore
cover every sufficiently small radius. The first crossing places a
prescribed radius in one of them. Diagonalize finite degree, arity,
window and derivative requests, requiring
$(\alpha_0+1)\log M_{n_i}\le(\log\widehat\beta_{n_i})^{48}$.
The overshoot in $k(n)$ is negligible compared with
$\widehat\beta_n$, as is the logarithm of the radius multiplier.
The running law in Chapter~\ref{ch10:trajectory} then gives
\eqref{ch17:running-radius}. This proves radius coverage. It remains to verify the uniform
comparison estimates used in these windows.

% !TEX root = ../monograph.tex
\subsection{Intervals of scales crossing changes of stage}
\label{ch17:variable-windows}

The finite comparison in Proposition~\ref{ch11:n-mark-comparison}
extends to the intervals of radii just described. We give the estimates
because requiring every interval of scales to lie inside a plateau would leave
gaps in radius coverage.

Fix the maximum scaling dimension of the insertions, their number, and
the orders of the derivatives to be estimated. In this paragraph
$C_j$ includes the linear column constants with insertions, the curvature-response
row constants and the finite polynomial-source propagation constants
for this request, all fixed \emph{before} the block ratio of stage $j$.
Write $\mathscr L_j$ for the ratio chosen at stage $j$, so
$L_i=\mathscr L_{J_i}$ with $J_i=J(i+1)$ in the scale convention
of Chapter~\ref{ch10:trajectory}. The ratio choices may require
\[
 \varepsilon_j:=\frac{\log(2+C_j)}{\log\mathscr L_j}\longrightarrow0.
\]
For example one may impose $\varepsilon_j\le j^{-2}$.
Constants depending on the already selected ratio, such as nonlinear
jet bounds of the critical configuration, are not included here. They enter the local
small errors and endpoint constants, and are handled by sufficiently large lower bounds on the reference coupling.
These local errors are summed after transfer; their ratio-dependent
constants are never multiplied once per step. Put $H_n=\exp((\log\log\widehat\beta_n)^2)$. Only a bounded number
of endpoint factors is multiplied, giving $H_n^{o(1)}$.
This distinction avoids a circular ratio choice.

For $b$ in a fixed compact interval of sufficiently large positive
numbers, move backwards from scale $n$, adding $\log L_i$ at each step.
Stop when this sum first reaches $(b+2)\log H_n$. Starting from that
scale, continue until the additional sum first reaches $10\log B_n$. Here $B_n$ is the radius of the restarted representation;
$\log B_n=O(\log\log\widehat\beta_n)$.
The schedule bound $L_i\le\log\widehat\beta_i$ makes each
excess over the target $O(\log\log\widehat\beta_n)=o(\log H_n)$.
The entire interval of scales has logarithmic length $O(\log H_n)$, and the
number of steps is $O(\log H_n)$ since every $L_i\ge3$.
It is therefore shorter than $(\log\log\widehat\beta_n)^4$,
as required by Proposition~\ref{ch11:restart-family}.
All its scale indices tend to infinity. The increment estimate is a
constant times the logarithmic block length plus the
local drift error controlled by the lower bound on the reference coupling. Summing it over this interval gives
\[
 \max_i|\beta_i/\beta_n-1|
 \le \frac{H_n^{o(1)}(1+\log H_n)}{\widehat\beta_n}.
\]
In particular coupling changes are smaller than every fixed inverse
power of $H_n$, not just qualitatively $o(1)$. An interval of scales may cross any number of stages; no bound by two
stages is needed.

For the linear propagation constants the exact bound is
\begin{equation}\label{ch17:window-product}
 \log\prod_{i\text{ in window}}(2+C_{J_i})
 \le\sup_{i\text{ in window}}\varepsilon_{J_i}
                \sum_{i\text{ in window}}\log L_i
 =o(\log H_n).
\end{equation}
Thus they contribute a total factor $H_n^{o(1)}$. In particular the differentiated
curvature-response estimate in Chapter~\ref{ch11:insertions} becomes
\[
 S^{r+2}\sum_y|\Delta^r(\Psi_{n-1}\cdots\Psi_\ell)(x,y)|
 \le 2^d\prod_{i=\ell}^{n-1}R_{J_i}=H_n^{o(1)},
 \qquad S=a_\ell/a_n.
\]
The covariance decomposition telescopes for these unequal
curvature responses just as it does for equal ratios: each conditional covariance
is conjugated by the product of the subsequent curvature responses.
It uses no replacement by a common stage's quadratic form.
The finite-cutoff covariance estimate \eqref{ch11:finite-covariance}
is uniform over the blocking sequence, in particular over the unequal ratios $L_i$ of the interval of scales (Section~\ref{ch11:window-covariance}). Its stage-dependent constants,
protected-path errors, ratio-dependent jet errors and finite incidence
costs are made smaller than the requested powers of $H_n$ by choosing
sufficiently large lower bounds on the reference coupling. Sums over $O(\log H_n)$ steps lose only a factor $H_n^{o(1)}$.
At each step the linear map has its norm bound fixed before the block
ratio is chosen. The remaining nonlinear term is small once the lower bound
on the reference coupling is sufficiently large. Constants depending on the chosen ratio
therefore enter the sum of local errors. The regular and large-field source recurrences therefore retain their
strict negative powers after multiplication by the growth factor
\eqref{ch17:window-product}.

One can now repeat the proof of
Proposition~\ref{ch11:n-mark-comparison} with unequal ratios.
The initial Taylor order and support margins are chosen for the
largest $b$ and the fixed finite request. The entire incoming
polynomial in the source variables is retained. The quadratic covariance
and the polynomial $\sum_IG_Iz_I$ of free moments are common to all subsets;
the density without insertions, and hence the partition function, is unchanged.
Absorbing the larger source weight still uses
$\omega=H_n^{1/(8N_*)}$ and $I=\log H_n/(16N_*)$.
Every additional factor lost in transfer is $H_n^{o(1)}$; the margins in the
$H_n^{-3/2}$ moment error, $H_n^{-31/32}$ connected error and
$H_n^{-1}$ bounds of Lemma~\ref{ch11:restore-regular} absorb it. The chessboard bound \eqref{ch11:n-mark-root} is
therefore at most $H_n^{-c_*}$ for a fixed $c_*>0$ depending only on
the finite arity, for example $c_*=1/(32N_*)$ after increasing the lower bound on the reference coupling.
Its cell side is at most $a_nH_n^{b+5}$.
The lower bound on the partition function is imposed at the stage of the coarsest scale of the interval of scales,
with its constant at most $I$; it is not multiplied once per step.

The common localization map $T$ depends only on the insertion
and its original self covariance. It is independent of the middle
and output stopping scales. Restart on the longest interval of scales and restrict
that restarted history to the shorter ones. This gives the same $T$ for all
$b$ in the compact interval. The finite number of coefficient bounds
above is uniform in $b$: changing $b$ only changes stopping integers
and the scaling dilation $W_{H_n^b}$, whose exponents lie in a
fixed compact interval. All upper bounds used the longest interval of scales and
the largest dilation factor. No compactness over arbitrary incoming
activity arrays is invoked.

Apply this construction to
$b=\alpha_j+u-\alpha_0$, $j=0,1,2$, and
$u\in[\alpha_0,\alpha_0+1]$.
At its own normalization a single insertion has bounded chessboard seminorm.
Changing from $\alpha_j$ to $\alpha_0$ contributes a factor at most
$CH_n^{(\alpha_0-\alpha_j)D}$ by the exact scaling dilation on
the common localization map. These are precisely the size bounds used
in the two estimates above for the coincident-coordinate strips and the grid. Their exponents are
strictly negative, so the harmless factors $H_n^{o(1)}$ do not change
the conclusion. The estimates with constants independent of the number of reflected copies give the same pair bounds
in every phase. The existing argument using separated correlation pairings then gives
fixed derivatives and shell fluxes with the same diagonal subsequence and the same comparison isomorphisms $U_{i,d}$.
Thus the prescribed-radius comparison used in this chapter follows
from the construction with insertions, including intervals of scales that cross
changes of stage.


\section{Tensor coordinates and the regulated current}
\label{ch17:current}

Use the free scalar and tensor
\[
 \mathcal O=\tfrac14\sum_{a,b,\alpha}{:}(F_{ab}^\alpha)^2{:},
 \quad L=T^{\rm Max},\quad D_{ab}=\delta_{ab}\mathcal O,
 \quad X_{ab}=(1-\delta_{ab})\sum_{c,\alpha}{:}F_{ac}^\alpha F_{bc}^\alpha{:}.
\]
The coordinate sum in $\mathcal O$ is over ordered pairs. The symbol
$X$ here denotes the off-diagonal tensor, not a polymer support.
These three tensors span the centered degree-four symmetric
hypercubic tensors. Before imposing the Maxwell equation, their
vector divergences have coordinates in $E=\RR^3$, represented by
\[
 \sum_c F_{bc}\partial_bF_{bc},\qquad
 \sum_{a,c}F_{ac}\partial_bF_{ac},\qquad
 \sum_{a,c}F_{bc}\partial_aF_{ac},
\]
with colour contraction and Wick ordering understood. The free
quotient $s:E\to V$ has kernel spanned by $v=(0,0,1)$, and $V$
is the two-dimensional vector space of Section~\ref{ch14:stress-line}.
For a tensor with diagonal entries
$z_0\sum_cF_{ac}^2+z_1\sum_{b,c}F_{bc}^2$ and off-diagonal entries
$z_2\sum_cF_{ac}F_{bc}$, Bianchi gives divergence coordinates
$(2(z_0-z_2),2z_1+z_2/2,z_2)$. Hence
\begin{equation}
 \rho(x,y,z)=-2x-8y,\qquad
 \rho(\operatorname{div}_{\rm raw}B)\mathcal O=-\operatorname{tr}B.
 \label{ch17:rho}
\end{equation}
In particular $\rho(v)=0$ and the values for $L,X,D$ are $0,0,-4$.

At a lattice site set
$f_{ab}=\operatorname{proj}_{\gLie}(U_{\square(a,b)}-
U_{\square(a,b)}^{-1})/2$ and start with
\[
 J_{ab}=\tfrac14\delta_{ab}\sum_{c,d}\langle f_{cd},f_{cd}\rangle
                  -\sum_c\langle f_{ac},f_{bc}\rangle,
 \qquad (Y_b^x)_e=\mathbf1_{e^-=x}f_{b,e^{\rm dir}}(x).
\]
Here $Y$ is a left link displacement. If $\mathsf S$ is the leading
path action without its coupling, define its current defect by
$\mathcal E_b=\sum_a\nabla_aJ_{ab}-D\mathsf S[Y_b]$, using
\[
 \nabla_a=\sum_{j=1}^4c_j(\mathsf s_a^j-\mathsf s_a^{-j}),
 \qquad(c_1,c_2,c_3,c_4)=(4/5,-1/5,4/105,-1/280).
\]
The operator $\mathsf s_a$ shifts by one site in direction $a$.
Initially the trace and antisymmetric part of $J$ vanish. During
transfer they must remain exact forward differences; smallness of their
values alone would not suffice for the trace rate below.

We fix the coefficient functionals in that calculation. For a local scalar
$F$, its Hessian Fourier kernel is
\[
 \widehat F_{c\alpha,c\alpha}(p,q)=
 \sum_{x,y}e^{i(p\cdot x+q\cdot y)}
 D^2F(0)[e_{x,c,\alpha},e_{y,c,\alpha}],
\]
where $e_{x,c,\alpha}$ is the link-coordinate basis vector. Fix one
colour $\alpha$. For a vector $F_b$ and a scalar action $F$ put
\begin{align*}
 \mathsf l(F)&=\operatorname{Re}\sum_{a\ne c,b}
      i[z^2t]\widehat{F_b}_{c\alpha,c\alpha}(ze_a,-ze_a+te_b),\\
 \mathsf q(F)&=|\Lambda|^{-1}\operatorname{Re}\sum_{a\ne c}
            [z^2]\widehat F_{c\alpha,c\alpha}(ze_a,-ze_a),\\
 \mathsf h(F)&=\mathsf l(F)-\tfrac12\operatorname{Re}\sum_{a\ne c}
            [z^2]\widehat{F_c}_{c\alpha,c\alpha}(ze_a,-ze_a).
\end{align*}
Brackets denote Taylor coefficients including reciprocal factorials.
The volume normalization in $\mathsf q$ is for the translation-invariant
action; $\mathsf l$ evaluates the anchored vector response. Define
$\mathfrak h(Y)=\mathsf h(\operatorname{div}_{\rm Haar}Y)$.
The second term in $\mathsf h$ corrects the lattice derivative.
The regulated current calculation uses
\[
 \Gamma_\mu=\tfrac12\log\det(H+dd^*+\mu^2I)
                 -\log\det(d^*d+\mu^2I),\quad
 R(\mu)=\mathsf q(\Gamma_\mu),\quad
 \mathcal D(\mu)=-\mu R'(\mu).
\]
Here $H$ is the left Hessian of $\mathsf S$, $d$ its infinitesimal
gauge insertion, and $\mathsf q$ is the coefficient functional just defined. Its normalization
counts twelve ordered transverse planes. The actual Ward calculation requires its Haar
coefficient $\mathfrak h(Y)$ to obey
\begin{equation}
 \mathfrak h(Y)=\mathcal D(\mu)+O(\mu)+o(1),\qquad
 R(\mu)=12\kappa_G\log\mu+O(1).
 \label{ch17:regulator-input}
\end{equation}
The $o(1)$ is uniform on each fixed interval $[\delta,1]$ as the
current defect tends to zero and the period tends to infinity;
the constant in $O(\mu)$ is independent of $\mu\le1$.
This identity concerns the actual gauge-fixed Hessian and transferred
current. The perturbative anomaly citation does not supply it.

The next section derives this regulated identity from explicit weighted
current bounds and convergence of the matched current jets. It keeps
the gauge-coordinate corrections and the lattice coefficient correction.
The following current-contraction section then verifies these hypotheses
for repeated stationary transfer.

% !TEX root = ../monograph.tex
% Fragment for Chapter17. Root owns integration and section numbering.
\section{Exact current transport and the scalar regulator}
\label{ch17:current-suppliers}

Several ingredients of the trace transfer can be established before its
quantitative current estimates. We give them separately, since they use
the actual difference operator and the actual path action. They do not
identify a lattice coefficient by comparison with a separately chosen
continuum action.

\subsection{Difference-preserving block averages}

Let $L$ be a block ratio and $s_a$ a unit fine translation. Write
$\delta_a=s_a-1$, $\delta_a^{\rm c}=s_a^L-1$, and let
$\nabla_a=\sum_{j=1}^4c_j(s_a^j-s_a^{-j})$ be the eighth-order
centered difference of Section~\ref{ch17:current}. Write
$\nabla_a^{\rm c}=\sum_{j=1}^4c_j(s_a^{Lj}-s_a^{-Lj})$ for the
corresponding coarse difference, acting on the fine lattice.
The coefficients satisfy
$2\sum_jjc_j=1$. Define finite Laurent polynomials
\begin{align*}
 M_a&=L^{-1}\sum_{v=0}^{L-1}s_a^v,&
 G_a&=\sum_{j=1}^4c_j\sum_{v=-j}^{j-1}s_a^v,\\
 N_a&=G_a(s_a^L)M_a,&
 K_0&=\prod_aM_aN_a,\\
 K_a&=G_aM_a\prod_{c\ne a}M_cN_c,&
 K_a^{\rm e}&=N_a\prod_{c\ne a}M_cN_c.
\end{align*}
All these operators act on the fine anchor of a local function before
stationary substitution. If $\mathcal I F=F(U_{\rm stat}(V))$, coarse
translation of $\mathcal IF$ is $\mathcal I(s_a^LF)$; fine unit shifts
are not silently identified with coarse shifts.

\begin{lemma}[Exact block identities]
\label{ch17:block-identities}
There is a finite Laurent polynomial $K_a^{\rm t}$ such that
\begin{align}
 K_0\nabla_a&=L^{-1}\nabla_a^{\rm c}K_a,&
 K_0\delta_a&=L^{-1}\delta_a^{\rm c}K_a^{\rm e},&
 K_a-K_0&=\delta_a^{\rm c}K_a^{\rm t}.
 \label{ch17:block-algebra}
\end{align}
The kernels have support of radius $CL$ and absolute coefficient sums
bounded by a numerical constant independent of $L$; their masses are one
except that no mass is prescribed for $K_a^{\rm t}$.
\end{lemma}
\begin{proof}
The identities $G_a\delta_a=\nabla_a$ and
$M_a\delta_a=L^{-1}\delta_a^{\rm c}$ prove the first two assertions
by commutation of Laurent polynomials. For the third, work in the
quotient ring with $s_a^L=1$. There $M_a^2=M_a$,
$M_aG_a=G_a(1)M_a=M_a$ and $G_a(s_a^L)=G_a(1)=1$.
Thus $M_a(G_a-G_a(s_a^L)M_a)$ is zero in that quotient, which is
precisely divisibility of $K_a-K_0$ by $s_a^L-1$.
The original kernels have finite products of a fixed number of bounded
coefficient-sum factors, each of range $CL$. For the quotient, group
coefficients in each residue class modulo $L$. Its numerator has zero
sum in each class, and only a bounded number of nonzero entries in that
class. Division by $s_a^L-1$ is a partial sum along this finite sequence,
so increases its absolute sum by only a bounded factor. This also proves
the support bound for $K_a^{\rm t}$.
\end{proof}

For a fine current with $\operatorname{tr}J=\sum_a\delta_aq_a$, set
\[
 J'_{ab}=L^4\mathcal I(K_aJ_{ab}),\qquad
 q'_a=L^3\mathcal I(K_a^{\rm e}q_a)
                       +L^4\mathcal I(K_a^{\rm t}J_{aa}).
\]
Equation~\eqref{ch17:block-algebra} proves exactly
$\operatorname{tr}J'=\sum_a\delta_a^{\rm c}q'_a$.
Indeed split $K_a=K_0+\delta_a^{\rm c}K_a^{\rm t}$ and then apply
the second identity to $K_0\delta_aq_a$.
Similarly, if $J_{ab}-J_{ba}=\sum_c\delta_ct_{abc}$, the transferred
antisymmetric part is the sum of the transferred old differences and
\[
 \delta_a^{\rm c} L^4\mathcal I(K_a^{\rm t}J_{ab})
 -\delta_b^{\rm c} L^4\mathcal I(K_b^{\rm t}J_{ba}).
\]
In particular symmetry may fail pointwise after blocking, but its failure
and the trace remain exact differences. This conclusion is independent
of estimates on the current's Taylor coefficients.

\subsection{What scalar mass regularization changes}

Consider the same admissible path action used in Chapter~\ref{ch07:matching},
with its left Hessian $H$ and infinitesimal gauge insertion $d$.
At the zero field the completed vector operator is
$H+dd^*=|k|^2I+O(|k|^3)$, and the ghost operator is
$d^*d=|k|^2I+O(|k|^4)$, with the global lattice ellipticity stated there.
Define $\Gamma_\mu$ with the scalar mass $\mu^2I$ as above.
The normalized quadratic response of $\Gamma_\mu$ has the same logarithmic
coefficient as the regulated response in Chapter~\ref{ch07:matching}.
This fact concerns its coefficient integral, before any relation to the
Haar divergence of a current.

To see this directly, retain the actual vector and ghost first and second
background insertions $b,c$ and replace only their denominators by the
scalar-mass denominators. In each channel the logarithmic Hessian is
\[
 \operatorname{tr}\bigl(a_\mu(k)^{-1}c(-p,p;k)
 -a_\mu(k)^{-1}b(-p;k+p)a_\mu(k+p)^{-1}b(p;k)\bigr).
\]
The scalar mass has no background insertion. The actual local Taylor
expansion and vertex estimates of Chapter~\ref{ch07:matching} therefore
apply without the regulator-vertex terms. With $T=|k|+\mu$, the
inverse and its $r$ momentum derivatives cost $CT^{-2-r}$, for every
fixed required $r$. Differentiate $Ra=I$ repeatedly to obtain this bound;
the denominator floor is now simply $a_\mu\ge cT^2I$.
Two external derivatives of the displayed expression cost at most
$CT^{-4}$. For $\mu\le|k|\le1$, comparison with the massless continuum
vertices by the resolvent identity costs
\[
 C\bigl(|k|^{-3}+\mu^2|k|^{-6}\bigr).
\]
The first term is the actual lattice Taylor remainder, and the second
is the scalar-mass denominator difference. Both have bounded integrals
in four dimensions, uniformly in $\mu$. The region $|k|<\mu$ also
has bounded integral, since its volume is $O(\mu^4)$ and its integrand
is $O(\mu^{-4})$. Away from zero the integral is uniformly bounded.
Thus the continuum homogeneous term alone determines the logarithm.
With the twelve-plane coefficient convention of this chapter this gives
\begin{equation}
 R(\mu)=12\kappa_G\log\mu+O(1),\qquad0<\mu\le1.
 \label{ch17:scalar-log-supplier}
\end{equation}
This is the same one-loop normalization already calculated using the
actual action in Chapter~\ref{ch07:matching}. At fixed $\mu>0$ the
integrand is periodic analytic on a common smaller strip, so its torus
sums converge to this integral; no uniform rate as $\mu\downarrow0$
is required. This proof does not differentiate the bounded remainder.

\subsection{Exact integration by parts and the remaining current estimate}

For a smooth vector field $Y$ on a finite product of compact link groups,
Haar integration gives the exact identity
\begin{equation}
 \int\bigl[\operatorname{div}_{\rm Haar}Y
                  -D\mathcal A[Y]\bigr]e^{-\mathcal A}\,dU=0.
 \label{ch17:finite-haar}
\end{equation}
It follows by integrating the divergence of $e^{-\mathcal A}Y$.
With a spectator $F$ the right side becomes
$-\int DF[Y]e^{-\mathcal A}\,dU$. In particular it is zero when
$Y$ acts only on links disjoint from those of $F$. Smooth truncations
can be inserted before this identity is used; their differentiated terms
must be retained as marked remainders.

There is also an exact cancellation available for the inherited quadratic
loop term. Chapter~\ref{ch07:matching} represents its Hessian as a sum
with two lattice differences on each leg. Consequently it annihilates
constant and affine polynomial link fields. The degree-three external
momentum coefficient in its variation by the leading current pairs a
constant leg with a quadratic leg, or two affine legs; both vanish.
This proves the zero trace-functional coefficient of that quadratic
loop contribution. It does not bound the nonlinear current remainder.

The missing quantitative step is now quite specific. One needs the
actual stationary pullback of the current defect, with the gauge-covariant block average
frame correction, in the exponentially weighted order-two, -three and
-four coefficient norms. In particular the order-two pullback must
annihilate its affine--affine and affine--quadratic polynomial moment conditions
\emph{before} estimating the remainder. The usual smooth transport bound
$\nabla_f^rW_a=O(L^{-1-r})$ alone does not assert these cancellations.
After gauge removal the relevant affine field costs $L^{-2}$ and its
quadratic term $L^{-3}$; multiplication by the current rescaling $L^5$
would leave the unacceptable factors $L$ and $1$ from their two products.
Only removal of these exact terms leaves the $L^{-1}$ remainder used
in the proposed contraction. Fourth-order action matching must therefore
be combined with an explicit identity for the current's matched jets,
not merely with a bound on the action's matching error.

The group-coordinate identity needed for this comparison is proved next.
Its proof is independent of the contraction of the actual current
sequence, established in Section~\ref{ch17:quadratic-current}.

\subsection{The group-coordinate contact identity}
\label{ch17:contact-section}

Fix a finite lattice and a background $U$ in the product of link groups.
Use the invariant inner product on its link Lie algebra and the chart
$z\mapsto e^zU$ (edge by edge). Let $S$ be a smooth gauge-invariant action,
$H=D_z^2S(e^zU)|_{z=0}$ its symmetric left-coordinate Hessian, and
$d$ the infinitesimal gauge map at $U$. The displacement $Y$ is assumed
gauge covariant, so its local flow
$\Psi_s(U)=e^{sY(U)}U$ commutes with gauge transformations. It need not
be the exact time-$s$ flow; all following identities use only its
first derivative at zero.

Define the chart velocity and its linear part by
\[
 \mathcal Y_U(z)=\left.\partial_s\right|_0
    \log\bigl(\Psi_s(e^zU)\Psi_s(U)^{-1}\bigr),
 \qquad V=D_z\mathcal Y_U(0).
\]
In particular $\mathcal Y_U(0)=0$. This definition includes the
change of the left frame as the base point moves. Replacing $V$ by an
ordinary derivative of the coefficient function $Y$ would omit that
change. Let
\[
 J_U(v,z)=\left.\partial_s\partial_t\right|_{0}
       \log\bigl(\operatorname{gauge}_{e^{sv}}(e^{tz}U)U^{-1}\bigr),
 \qquad N_{i\alpha}=DS(U)[J_U(e_\alpha,e_i)].
\]
Indices $i$ run through link coordinates and $\alpha$ through gauge
coordinates, using the basis of Section~\ref{ch01:gauge} at each link and site. Differentiating gauge invariance first in $v$, then in $z$,
gives $Hd=-N$.

For $\mu>0$ at a background where the following inverses exist, write
\[
 R=(H+dd^*+\mu^2I)^{-1},\qquad
 G=(d^*d+\mu^2I)^{-1},
\]
\[
 \Gamma=\tfrac12\log\det(H+dd^*+\mu^2I)
                            -\log\det(d^*d+\mu^2I).
\]
Near the zero background the positive denominators of the preceding
subsection make these ordinary real analytic functions. Define a link
vector $Z$ by its action against an arbitrary covector $\ell$:
\begin{align}
 \ell(Z)={}&\tfrac12\sum_{i,j}R_{ij}\,
                   \ell\bigl(D_z^2\mathcal Y_U(0)[e_i,e_j]\bigr)\notag\\
 &-\sum_{i,\alpha}
       [Gd^*(V+V^*)R]_{\alpha i}\,\ell(J_U(e_\alpha,e_i)).
 \label{ch17:contact-Z}
\end{align}
All matrices and derivatives here belong to the same background and
coordinate chart.

\begin{proposition}[Exact contact identity]
\label{ch17:contact-identity}
With these definitions,
\begin{equation}
 \tfrac12\operatorname{tr}\bigl(RD^2(DS[Y])\bigr)
       -\operatorname{div}_{\rm Haar}Y
       =D\Gamma[Y]-\mu^2\operatorname{tr}(RV)+DS[Z].
 \label{ch17:contact}
\end{equation}
No stationarity of $U$ is assumed. In particular the $DS[Z]$ term
cannot be discarded before applying the relevant coefficient functional.
\end{proposition}
\begin{proof}
Let a dot denote differentiation along $\Psi_s(U)$, using the left
chart anew at that base point. The chain rule for
$S(\Psi_s(e^zU))$ gives
\[
 D^2(DS[Y])=\dot H+V^*H+HV+C,
 \qquad C_{ij}=DS[D_z^2\mathcal Y_U(0)[e_i,e_j]].
\]
This is an exact chart calculation, not a replacement of the group by
its tangent space. The differential of $\Psi_s$ in these moving charts
is $I+sV+o(s)$. Exponential charts have Haar density one at their origins,
so its Jacobian at the base point has derivative
$\operatorname{div}_{\rm Haar}Y=\operatorname{tr}V$.

Gauge equivariance of $\Psi_s$ gives $\dot d=Vd$: differentiate
$\Psi_s(\operatorname{gauge}_{e^v}U)
=\operatorname{gauge}_{e^v}\Psi_s(U)$ first in $v$ and then in $s$.
Consequently
\[
 D\Gamma[Y]=\tfrac12\operatorname{tr}(R\dot H)
       +\operatorname{tr}(dd^*RV)
       -\operatorname{tr}\bigl(dGd^*(V+V^*)\bigr).
\]
On the other hand the Hessian chain rule and
$RH=I-Rdd^*-\mu^2R$ yield
\[
 \tfrac12\operatorname{tr}(RD^2(DS[Y]))-\operatorname{tr}V
 =\tfrac12\operatorname{tr}(R\dot H)
       -\operatorname{tr}(Rdd^*V)-\mu^2\operatorname{tr}(RV)
       +\tfrac12\operatorname{tr}(RC).
\]
Finally $Hd=-N$ implies the exact resolvent identity
\[
                  Rd=dG+RNG.
\]
Subtract the two preceding displays. The leftover gauge terms are
\[
 -\operatorname{tr}\bigl((Rdd^*+dd^*R-2dGd^*)V\bigr)
       =-\operatorname{tr}\bigl(Gd^*(V+V^*)RN\bigr).
\]
Inserting the definition of $N$ and the term $\frac12\operatorname{tr}(RC)$
identifies their sum with $DS[Z]$ from \eqref{ch17:contact-Z}.
This proves \eqref{ch17:contact}.
\end{proof}

\subsection{Cancellation of the action-gradient terms}

The functional $\mathsf l$ above extracts the coefficient of $z^2t$
from the diagonal polarization pair at momenta
$p=ze_a$, $q=-ze_a+te_b$, and sums over $a\ne c$ and $b$.
We check why it kills the action variations in
\eqref{ch17:contact}. Work at fixed positive $\mu$ so all kernels
are analytic in momentum near zero. A gauge-invariant, translation- and
hypercubic-invariant scalar action $F$ has quadratic kernel
$c_F(|p|^2I-p\otimes p)+O(|p|^3)$ by the Ward calculation in
Chapter~\ref{ch07:matching}. This applies to both $S$ and $\Gamma$.
Let $Z_b$ be any anchored gauge-covariant displacement, including $Y_b$
or the displacement \eqref{ch17:contact-Z}. Simplicity of the gauge
algebra implies $Z_b(0)=0$. The linearized covariance identity says
$DZ_b(0)d_0=0$.

Use $i$ for its output link direction and $c$ for its input polarization.
Its joint momentum symbol, with output dual variable $p$ and input
variable $q$, therefore has expansion
\[
 (DZ_b)_{ic}(p,q)=\sum_j q_j A_{b,ic}^{\,j}+O((|p|+|q|)^2),
 \qquad A_{b,ic}^{\,j}=-A_{b,ij}^{\,c}.
\]
The factors $i$ from Fourier differentiation may be absorbed into $A$.
There is no term depending only on $p$: the identity obtained by
multiplying by the gradient symbol in $q$ forces it to vanish.
For $a\ne c$, the coefficient of $z^2t$ in the quadratic variation
$DF[Z_b]$ is, apart from the common factor $c_F$,
\[
       A_{b,cc}^{\,b}-2\delta_{ab}A_{b,cc}^{\,a}
                              +\delta_{bc}A_{b,ac}^{\,a}.
\]
This follows by inserting the two kinetic kernels on its two legs.
The first sum over $a\ne c$ gives
$3\sum_{b,c}A_{b,cc}^{\,b}$. The second gives minus twice that sum,
since $A_{c,cc}^{\,c}=0$. The last gives minus the same sum after
using $A_{c,ac}^{\,a}=-A_{c,aa}^{\,c}$ and exchanging $a,c$.
Thus
\begin{equation}
                  \mathsf l(DF[Z])=0.
 \label{ch17:gradient-cancellation}
\end{equation}
Higher momentum orders cannot contribute to the coefficient of degree
three. Gauge covariance of \eqref{ch17:contact-Z} follows by transforming
its chart tensors, inverse operators and contracted indices together.
The mass $\mu^2I$ is unchanged by these orthogonal gauge frame changes.
The calculation therefore applies to both action-gradient terms of the
exact contact identity. It is a finite-momentum cancellation, not a
statement that the corresponding functions vanish pointwise.

\subsection{Power counting the regulator remainder}

The remaining estimate has a transparent sufficient numerator condition.
Let $W^{(1)},W^{(0)}$ be the vector and ghost remainder operators obtained
by subtracting the leading curvature displacement from $V$; in the
notation of the current calculation these are
\[
 W_b^{(1)}=\operatorname{sym}(V_b-dM_b-D_b)+\tfrac12T_b^{(1)},
 \qquad W_b^{(0)}=\operatorname{sym}(M_bd)-\tfrac12T_b^{(0)}.
\]
Here $\chi$ is the indicator of the anchored lattice site,
$(M_bh)(x)=\chi(x)h_b(x)$,
$(D_bh)_a(x)=\delta_a\chi(x)h_b(x)$, and
$(T_b^{(r)}h)(x)=\delta_b\chi(x)h(x)$ on $r$-forms.
The full lattice shifts and chart tensors must be retained when
extracting their coefficient; their continuum symbols only identify
the leading terms that cancel.

For that leading calculation the gauge derivative is $d=-\partial$
and the curvature displacement is
$(V_bh)_a=\chi(\partial_bh_a-\partial_ah_b)$.
Then $dM_b+D_b=-\chi\partial_ah_b$ and
$\operatorname{sym}(\chi\partial_b)=-\frac12\partial_b\chi$.
Thus both $W^{(r)}$ have zero first-order principal symbol. Covariantizing
these expressions gives the same cancellation at constant background
to first field order: the skew-adjoint multiplication by
$\operatorname{ad}A_b$ has zero symmetric part, and the remaining
$\operatorname{ad}A_a h_b$ term cancels the corresponding term of $dM_b$.
Accordingly their zero-, one- and two-field symbols have orders
$2,1,0$, respectively. A lattice/current sequence having the stated
matched free jets inherits precisely these vanishings in its limit.
Weighted coefficient bounds control the Taylor remainders.

Here is the resulting estimate without suppressed differentiation orders.
Set $z=|q|+\mu$. For $n\le2$ background derivatives and $r\le3$ external
momentum derivatives the inverse estimates are
$Cz^{-2-n-r}$. The remainder symbols have bounds
$Cz^{2-n-r}$. The quadratic-field coefficient of
$\operatorname{tr}(RW^{(1)})+\operatorname{tr}(GW^{(0)})$, after the
three momentum derivatives defining $\mathsf l$, is therefore bounded
by a sum of terms
\[
 Cz^{-2-n-r}\,z^{2-(2-n)-(3-r)}=Cz^{-5}.
\]
In four dimensions
$\int_{|q|\le1}(|q|+\mu)^{-5}\,dq\le C/\mu$.
The region away from zero is bounded. Hence the regulator factor
$\mu^2$ gives a remainder at most $C\mu$.
For approximate matched jets, first fix $\mu\ge\delta>0$; their
errors then tend to zero uniformly on that regulator interval.
Only afterwards send $\delta$ to zero. This is the order required
by the logarithmic mean-value argument.

\subsection{The corrected coefficient and order of limits}
\label{ch17:regulator-completion}

Retain the lattice correction in $\mathsf h$. For the same unit vector
$u\in\gLie$, set
\[
 \mathsf c(F)=\tfrac12\operatorname{Re}\sum_{a\ne c}[z^2]
 \widehat{F_c}_{cu,cu}(ze_a,-ze_a),
 \qquad \mathsf h=\mathsf l-\mathsf c.
\]
The quadratic variation $DF[Z]$ in the preceding calculation starts at
momentum degree three. Thus $\mathsf c(DF[Z])=0$, in addition to
\eqref{ch17:gradient-cancellation}. Both functionals therefore annihilate
the action-gradient terms of \eqref{ch17:contact}.

Suppose now that the actual current sequence has the following
properties at the fixed output level: uniform exponentially weighted
order-two through order-four coefficient bounds; convergence of the
needed low polynomial jets to the curvature current; the exact trace
and antisymmetry difference identities; and defects tending to zero
through order four in those weighted norms in
\[
                 DS[Y_b]\equiv_4\sum_a\nabla_aJ_{ab}-\mathcal E_b.
\]
The notation $\equiv_4$ refers to equality of field derivatives through
order four at the zero background. It suffices because the contact
expression uses two internal and two external field derivatives.
At fixed $\mu\ge\delta>0$, contracted defect errors tend to zero.
Translation covariance makes contraction by the inverse $R$ commute
with translating the anchor. The $\mathsf l$ coefficient of the
contracted current divergence selects $a=b$ at its sole linear
 total-momentum derivative. The sum is the contracted trace of $J$,
which has a further total-momentum zero because it is a forward
difference. Its coefficient is zero. For $\mathsf c$, the total
momentum is zero from the outset and the divergence vanishes directly.
The exact contact identity consequently yields
\[
 \mathfrak h(Y)=\mu^2\mathsf h\bigl(\operatorname{tr}(RV_b)\bigr)+o(1),
\]
uniformly on every fixed positive regulator interval. All operators here
are functions of the background before the two field derivatives are
extracted.

The definitions of $W^{(1)},W^{(0)}$ give an exact trace decomposition.
Using the same link and site bases induced by Section~\ref{ch01:gauge}, apply $Rd=dG+RNG$ to obtain
\begin{align*}
 \operatorname{tr}(RV_b)
  ={}&\operatorname{tr}\bigl(R(D_b-\tfrac12T_b^{(1)})\bigr)
      +\tfrac12\operatorname{tr}(GT_b^{(0)})+DS[Z_b']
      +\mathrm{Rem}_b,\\
 \mathrm{Rem}_b={}&\operatorname{tr}(RW_b^{(1)})
                                   +\operatorname{tr}(GW_b^{(0)}),\\
 \ell(Z_b')={}&\sum_{i,\alpha}(GM_bR)_{\alpha i}
                                       \ell(J_U(e_\alpha,e_i)).
\end{align*}
For example, the part containing $dM_b$ becomes
$\operatorname{tr}(GM_bd)+\operatorname{tr}(RNGM_b)$;
its symmetric ghost part is precisely
$\frac12\operatorname{tr}(GT_b^{(0)})+\operatorname{tr}(GW_b^{(0)})$,
and its remaining term is $DS[Z_b']$. This also fixes the sign of
the gauge correction.

For the three displayed principal terms the Fourier transform of
$\delta_b\chi$ supplies $e^{-it}-1=-it+O(t^2)$ at total momentum
$te_b$. With the factor $i$ in $\mathsf l$, the $D_b$ terms sum to
$\mathsf q(\operatorname{tr}R)$: the selected derivative has $a=b$
and the component trace then sums its diagonal entries. Each of the
four $T_b^{(1)}$ terms gives $\mathsf q(\operatorname{tr}R)$, and each
$T_b^{(0)}$ term gives $\mathsf q(\operatorname{tr}G)$.
The volume normalization is exactly the translation average of the
single anchor. Their total is therefore
\[
           -\mathsf q(\operatorname{tr}R)
                              +2\mathsf q(\operatorname{tr}G).
\]
Their $\mathsf c$ coefficients are zero, since $\delta_b\chi$ has
zero symbol at total momentum zero. Multiplication by $\mu^2$ gives
$-\mu\partial_\mu\mathsf q(\Gamma)$ by differentiation of the two
determinants. Thus, with the lattice correction fully retained,
\begin{equation}
 \mathfrak h(Y)=\mathcal D_\Lambda(\mu)
                  +\mu^2\mathsf h(\mathrm{Rem})+o(1).
 \label{ch17:corrected-regulated-ward}
\end{equation}
The action-gradient term $DS[Z_b']$ is annihilated by both parts of
$\mathsf h$ as proved above.

The $\mathsf l$ remainder was bounded by $C\mu$ after multiplication
by $\mu^2$. The correction $\mathsf c$ uses only two external momentum
derivatives, so the same product rule costs $Cz^{-4}$. Its integral
is at most $C(1+|\log\mu|)$ in four dimensions, and
$\mu^2(1+|\log\mu|)\le C\mu$ for $0<\mu\le1$.
Consequently \eqref{ch17:corrected-regulated-ward} proves the regulator
estimate, including the anchored correction, from the stated current
hypotheses.

Finally combine it with \eqref{ch17:scalar-log-supplier}. On a fixed
interval $[u-A,u]$ in the variable $\log\mu$, the bounded error in
$R(\mu)-12\kappa_G\log\mu$ implies that some regulator has
$|\mathcal D_\Lambda(\mu)+12\kappa_G|\le(2C+o(1))/A$ after the
period is large. The remainder is at most $Ce^u+o(1)$ on that interval.
First take the current and period limits there, next $u\to-\infty$,
and finally $A\to\infty$. Since $\mathfrak h(Y)$ is independent of
$\mu$, this proves
\[
                         \mathfrak h(Y)\longrightarrow-12\kappa_G.
\]
The group-coordinate contact identity, coefficient cancellation and
regulator estimate are now established. Their application still requires
the actual transferred current sequence to satisfy the weighted bounds
and matched-jet convergence stated at the start of this subsection;
those conditions are not consequences of the contact algebra.

% !TEX root = ../monograph.tex
% Current defect and uniform transfer bounds.
\section{Current conservation and uniform transfer bounds}
\label{ch17:quadratic-current}

The potentially dangerous quadratic defect has an exact cancellation
that should be imposed before estimating its norm. We prove preservation
of that cancellation for the linearization of the critical configuration. This removes
an additive interpretation of the order-two matching error. We then control the nonlinear
jets of the critical configuration and the current itself.

Fix a stage and an admissible action. Let $\mathcal E_b$ be a
gauge-invariant current defect with a base point, with zero constant and linear
field coefficients. Its Hessian at zero has an exponentially local
joint Fourier symbol $E_b(p,q)$, with two independent field momenta.
Gauge invariance gives a Ward identity on each leg. We say that its
quadratic defect has order at least four when
\begin{equation}
                    E_b(p,q)=O((|p|+|q|)^4)
                              \quad\hbox{near }(0,0).
 \label{ch17:defect-four}
\end{equation}
Equivalently, every Taylor coefficient of $E_b$ of total degree less
than four in $(p,q)$ vanishes. This order of vanishing in momentum
implies the cancellations for pairs of affine fields and for an affine
field paired with a quadratic field that are used in the current estimate.

For the initial Wilson current, \eqref{ch17:defect-four} follows from
the leading continuum identity. Its linearized curvature is $dA$;
the leading quadratic action is the curvature square; and the centered
difference $\nabla$ has linear symbol $ip$. The quadratic Taylor
coefficient of
$\sum_a\nabla_aJ_{ab}-DS[Y_b]$ therefore has leading spatial
momentum degree three. At that degree it is exactly
\[
 \sum_a\partial_a\left(\tfrac14\delta_{ab}\sum_{c,d}\|F_{cd}\|_{\gLie}^2
                         -\sum_c \langle F_{ac},F_{bc}\rangle_{\gLie}\right)
       +\sum_{a,c} \langle F_{bc},\partial_aF_{ac}\rangle_{\gLie}=0.
\]
The identity uses the linear Bianchi equation. Lower spatial degrees vanish because each curvature
contains a difference and the action variation contains one more.
All remaining lattice Taylor terms have degree at least four.
The same reasoning applies to any initial path action with the
normalized Taylor polynomial of scaling dimension four already established in
Chapter~\ref{ch07:matching}.

\subsection{Preservation of the order of vanishing}

The first derivative of the nonlinear critical-configuration map, from coarse
to fine link variables with block ratio $L$, is the forest representative
$W_F$ of \eqref{ch05:stationary-tangent}; its higher derivatives
cannot enter the quadratic pullback because the defect has zero first
field derivative. Since $W_F-W_a$ is a fine gradient, the Ward identities
of $E_b$ allow us to compute with the smooth representative $W=W_a$.

For the exact Taylor-order statement we pass once more to a gauge-equivalent representative, the linearization $W_s$ built with the linear averaging operator $M$ (Section~\ref{ch05:actual-tangent}), and test Hessians on polynomial fields.\mgubi{Issue 8: polynomial fields live on $\ZZ^4$, the critical configuration of Chapter~\ref{ch05:gaussian} on a torus. One sentence is missing: the Hessian on the torus is the periodization of the one on $\ZZ^4$. Vanishing on a constant leg (used below) is false on a torus. Counterparts: existence/\allowbreak quadratic-pole-transport.\allowbreak tex; v2 I-section.\allowbreak tex.} By \eqref{ch05:corrected-transport}, $W_a-W_s=-d_{0,f}ERW_s$; since $E$ is bounded and $RW_s$ has a bounded coarse kernel, each column of this difference is the gradient of a bounded fine potential. Let $F$ be a gauge-invariant local function of the fine field with base point $o$ and zero first field derivative, and write $B(a,a')=D^2F(0)[a,a']$, whose kernel decays exponentially in the distances of both legs from $o$.
The Ward identity on each leg, $B(d_{0,f}\phi,a')=B(a',d_{0,f}\phi)=0$, holds for finitely supported $\phi$ and, by this decay, also for bounded $\phi$ and fields $a'$ of polynomial growth: replace $\phi$ by its restriction to a large ball about $o$. Call a one-form with values in $\gLie$ a polynomial field of degree $r$ if its components in a basis of $\gLie$ are polynomials of degree at most $r$ in the coordinates. $B$ converges absolutely on pairs of polynomial fields, and the Taylor coefficient of multidegree $(\alpha,\alpha')$ at zero of the joint symbol of $B$, expanded about $o$, is a nonzero multiple of the value of $B$ on the fields $(x-o)^\alpha e$ and $(x'-o)^{\alpha'}e'$ with $e,e'$ constant. The symbol being analytic, \eqref{ch17:defect-four} for $F$ is therefore equivalent to
\begin{equation}
 B(P,P')=0\qquad\hbox{for polynomial fields $P,P'$ of degrees $r,r'$ with }r+r'\le3.
 \label{ch17:polynomial-four}
\end{equation}
This condition does not depend on the center of the coordinates, since a translate of a polynomial has the same degree. Hence it holds for every fine translate of $E_b$, and so for $F=K_0E_b$, where $K_0$ is the finite combination of fine translates of the base point from Lemma~\ref{ch17:block-identities}; this $F$ is again gauge invariant with zero first field derivative, and its kernel decays exponentially because $K_0$ has range $CL$.
Let $G=\mathcal IF$ be its composition with the critical configuration. For finitely supported coarse fields $a,a'$ the preceding paragraph gives $D^2G(0)[a,a']=B(W_Fa,W_Fa')=B(W_aa,W_aa')$, where the potential of the fine gradient $W_Fa-W_aa$ is bounded because it is a sum of values of $W_aa$ along the spanning tree of one block; and the Ward identity applied to the bounded potentials above gives $D^2G(0)[a,a']=B(W_sa,W_sa')$. The kernel of $W_s$ is $O(L^{-1})$ with exponential coarse decay (Section~\ref{ch05:transports}).
Hence the coarse Hessian kernel of $G$ decays exponentially in the distances of both legs from the coarse block of $o$, its joint symbol is analytic, and for coarse polynomial fields $P,P'$ absolute convergence gives $D^2G(0)[P,P']=B(W_sP,W_sP')$. As for $F$, \eqref{ch17:defect-four} for $G$ is equivalent to the vanishing of these values when $r+r'\le3$.
By \eqref{ch17:polynomial-response}, whose derivation in the proof of Proposition~\ref{ch17:spatial-moments} uses only the formulas of Chapter~\ref{ch05:gaussian}, $W_s$ maps a coarse polynomial field $P$ of degree $r<j$ to a fine polynomial field of degree at most $r$; its homogeneous part of degree $r$ is $L^{-r-1}$ times that of $P$, read in fine coordinates. Here $j$ is the smoothing order of the averaging weight, and $j\ge5$, since \eqref{ch05:actual-bridge} is used below with three fine differences.\mgubi{Issue 8: $j\ge5$ is stronger than the standing $j\ge4$ of Chapter~\ref{ch05:gaussian} and is not entered in the schedule of Chapter~\ref{ch10:trajectory}; nor are the requirements on the smoothing order in Section~\ref{ch11:marked-step} and in the proof of Proposition~\ref{ch17:spatial-moments}.}
If $r+r'\le3$, the two images therefore have degrees with sum at most three, and \eqref{ch17:polynomial-four} for $F$ gives $D^2G(0)[P,P']=0$. Thus $G$, and the pulled-back defect $L^5G$, satisfy \eqref{ch17:defect-four} on the coarse lattice.
The same computation gives the exact rescaling of the first coefficients that need not vanish. Suppose that $B$ vanishes on all pairs of polynomial fields with $r+r'<n$, and let $P,P'$ be homogeneous of degrees $r,r'$ with $r+r'=n<j$. Then the parts of $W_sP$ and $W_sP'$ of lower degree contribute nothing, and $D^2G(0)[P,P']=L^{-n-2}B(P,P')$, with $P,P'$ read in fine coordinates on the right.
No alias decomposition of the linearization and no projection off the fine gradient direction is used: fine gradients are removed by the Ward identity before any polynomial is inserted, and only the polynomial moments of $W_s$ enter.
No bound uniform in $L$, and no inverse of a kinetic symbol at a nonzero alias, is
needed for this exact Taylor-order statement.

The base-point average $K_0$, the average over translates of the base point, from
Lemma~\ref{ch17:block-identities}, is performed before composition with the critical
configuration. Its symbol has a double zero at every nonzero
alias: at least one coordinate of $l$ is nonzero, and the factor
$M_a^2G_a(s_a^L)$ has two factors of the average $M_a$ there.
The argument above does not use this double zero. Of $K_0$ it uses only
that it is a finite combination of fine translates of the base point; the
behaviour of the linearization enters only through \eqref{ch17:polynomial-response}.

Thus the order of vanishing is preserved by an exact computation. The
pulled-back Hessian, evaluated on pairs of coarse polynomial fields of total degree
at most three, equals the Hessian of $K_0\mathcal E_b$ evaluated on their images
under $W_s$. These images are fine polynomial fields of no larger total degree,
on which that Hessian vanishes by \eqref{ch17:polynomial-four}. No estimate,
no limit and no bound uniform in $L$ enters. Hence $L^5\mathcal I(K_0\mathcal E_b)$
satisfies \eqref{ch17:defect-four} on the coarse lattice. This is an
exact identity and holds on the same history.

In particular the quadratic jet is not merely close to having the
required affine cancellations: it has them exactly, provided the
transfer of the current uses the exact pulled-back defect
$\mathcal E'_b=L^5\mathcal I(K_0\mathcal E_b)$ through the relevant
field order. If an implementation redefines that defect by action
matching, equality of these quadratic jets must be retained; a bounded
additive matching error would not preserve this argument.

\subsection{A homogeneous bound for the quadratic channel}

The analytic division argument of Section~\ref{ch07:division} can be
applied to each of the two independent momentum legs of $E_b$.
The Ward identities produce one curvature difference on each leg.
Their homotopy divisions lower Taylor order by one each, so
\eqref{ch17:defect-four} leaves a coefficient of total order at least
two. Coordinatewise Taylor division in the remaining eight variables
factors two additional differences, distributed between the legs.
On a fixed smaller coefficient width these operations are bounded by
constants depending only on the stage and the decrease of the width. Thus $E_b$
is a sum of terms with either one and three differences, or two and
two differences, on its two field legs. The coefficient sum is bounded
by $C_j\|E_b\|$ in the original weighted norm.

Recall that $W=W_a$ is the smooth representative; the quadratic Ward
identities annihilate its gradient difference from $W_F$.
Its fine-difference bounds are
\[
                  \|\nabla_f^rW_a\|\le C_{j,r}L^{-1-r},
                                  \qquad r\le3.
\]
Consequently either allocation of four differences contributes at most
$C_jL^{-6}$ for the two legs. The base-point average $K_0$ has bounded
absolute coefficient sum, not a factor $L^4$; it is an average rather
than an unnormalized sum. Translated coefficient weights are absorbed
by a fixed narrower coarse width, exactly as in the inherited quadratic
action estimate of Chapter~\ref{ch07:matching}. Multiplication by the
current normalization $L^5$ proves
\begin{equation}
             \|D^2\mathcal E'_b(0)\|_{\rm out}
                  \le \frac{C_j}{L}\,
                            \|D^2\mathcal E_b(0)\|_{\rm in}.
 \label{ch17:quadratic-current-contraction}
\end{equation}
This is homogeneous in the defect. No constant inhomogeneous term has
been introduced. Dividing the coefficients by $(\log\widehat\beta)^2$
only adds the adjacent logarithmic ratio, which is bounded and tends
to one on a fixed-stage tail. The selected block ratio makes this
quadratic channel a strict contraction.

\subsection{Ratio-dependent bounds for the nonlinear terms}

The nonlinear jets of the critical configuration enter only below the diagonal in field
order. Their constants may depend on $L$, because the lower bound on the reference coupling is
chosen after the ratio. What must remain uniform before $L$ is the
coarse coefficient width on which these constants are measured.
We obtain that width from the gauge-fixed inverse of
Chapter~\ref{ch07:matching}.

Write $z(v)$ for the fluctuation coordinates of the critical configuration, and
$\mathscr S(v,z)$ for the same joint classical action as there.
Let $T u=(u,Dz(0)u)$, and let $\mathscr H$ be its Hessian in the fluctuation coordinates.
On the slice of the block axial gauge implicit differentiation gives
\begin{align}
 \mathscr H z_2[u,v]
    &=-D^3\mathscr S[T u,T v,\,\cdot\,],\notag\\
 \mathscr H z_3[u,v,w]
    &=-D^4\mathscr S[T u,T v,T w,\,\cdot\,]
      -\sum_{\{ij\}\{k\}}
             D^3\mathscr S[(0,z_2[u_i,u_j]),T u_k,\,\cdot\,].
 \label{ch17:stationary-higher-jets}
\end{align}
There are three terms in the sum. The right sides are fluctuation
covectors, with all mixed derivatives in the averaging operation and in the local gauge transformation included.

These formulas have period-independent bounds $C_{j,L}$ on one
fixed smaller coarse Fourier strip. To check the width statement in this zero-background computation, take the forest of
Section~\ref{ch05:actual-coordinates} to consist of the coarse translates of one reference tree. This fixes a free choice for the present computation only; the average over forests of Section~\ref{ch09:forest-average} is not used here, and independence of the forest is not claimed for the forest-dependent jets $z_2,z_3$ themselves. The bounds are used below only for pullbacks of gauge-invariant functions along the critical-configuration map; that these do not depend on the forest (the slices of the block axial gauge for two forests differ by a gauge transformation equal to one at the block roots) is used here and is not proved in this chapter.\mgubi{Issue 7: v3 has no such lemma. Lean has it near the zero configuration (YM.\allowbreak Work.\allowbreak F05.\allowbreak answer, from graph\_switch, fine\_switch, Forest.\allowbreak fold\_root; the proofs of the first two were not read). For the independent forests of Chapter~\ref{ch09:renormalization}, $\Pi_F$ does not commute with coarse translations.} Let $\Pi_F$ be projection onto the slice of the block axial gauge along the joint gauge
directions; it then commutes with coarse translations. At zero it is computed by finite sums along the block tree
and the local gauge transformation at the root; its coefficients and those of $\Pi_F^*$ have norm
bounded by a finite $C_{j,L}$, but bounded range in coarse blocks.
The dual $\Pi_F^*$ extends a slice covector by zero on gauge directions.
For such a covector the gauge-fixed inverse $\mathscr C^{-1}$ of
Chapter~\ref{ch07:matching} solves the Hessian equation modulo gauge.
Thus the slice inverse is
\[
                       \Pi_F\mathscr C^{-1}\Pi_F^*.
\]
The gauge-fixed inverse has a common positive coarse width chosen before
$L$; multiplication by the finite-range projections does not destroy it.
The path vertices and derivatives of the averaging operation also have a common coarse
width, with finite $C_{j,L}$ norms. Restrict the external strips by a
fixed factor so all intermediate momentum sums lie in this inverse
strip. Formula~\eqref{ch17:stationary-higher-jets} then proves the
asserted bounds for $z_2,z_3$. It does not assert bounds uniform in $L$
for their norms.

For every gauge-invariant function with a base point, evaluate this critical
configuration in a radial gauge at the base point. At first order, the gauge-corrected
link at fine displacement $x$ is a sum of the curvature of the linearization over
coordinate rectangles of total area at most $C(1+|x|)$. The smooth
curvature bound for $W_a$ gives
\[
 |W^{\rm rad}(x,y)|
 \le C_jL^{-2}(1+|x|)
                   e^{C|x|/L}e^{-c_jd(0,y)}.
\]
The polynomial factor and fine exponential are absorbed by the weighted
local coefficient norm. A common strip reduction handles the coarse
weight. Higher derivatives of this gauge change are finite products
along the same contours; the fixed-order factors are polynomial in contour
length. Their period-independent norms are therefore still finite
$C_{j,L}$ on the same smaller coarse width.

The raw higher coefficients $D^3\mathcal E,D^4\mathcal E$ need not
separately satisfy linear Ward identities. We use gauge invariance of
the complete function to make this radial substitution first; its
nonlinear derivatives multiply only lower field derivatives of the
defect.

Now differentiate the exact pullback $\mathcal E'=L^5\mathcal I(K_0\mathcal E)$.
In field order $r$ the term involving $D^r\mathcal E$ has only $r$
factors of the linearization of the critical configuration. For $r=3,4$, radial gauge bounds it by
$C_jL^{5-2r}\|D^r\mathcal E\|$. The quadratic diagonal has the
stronger estimate \eqref{ch17:quadratic-current-contraction}.
Every other chain-rule term involves $D^s\mathcal E$ with $s<r$
and higher jets of the critical configuration or of the gauge change. Hence, for $r=2,3,4$,
\begin{equation}
 e'_r\le a_{j,r}(L)e_r+\sum_{2\le s<r}C_{j,L,r,s}e_s,
 \qquad
 a_{j,2}(L),a_{j,3}(L)\le C_j/L,
 \quad a_{j,4}(L)\le C_j/L^3.
 \label{ch17:unnormalized-defect}
\end{equation}
Here $e_r$ is the weighted norm of the field derivative of order $r$.
The estimates are homogeneous in the defect, because every term in the
composition contains one of its derivatives. There is no action-independent
inhomogeneous term on the right.

Put $\bar e_{m,r}=e_{m,r}/(\log\widehat\beta_m)^r$. For a fine-to-coarse
step from $m+1$ to $m$, the coefficient below the diagonal becomes
\[
 C_{j,L,r,s}
       \frac{(\log\widehat\beta_{m+1})^s}
            {(\log\widehat\beta_m)^r},\qquad s<r.
\]
At fixed $j,L$ this tends to zero as the lower bound on the reference coupling increases.
The diagonal logarithmic ratio tends to one. Choose $L$ first so the
three diagonal coefficients are at most $1/8$, then choose the lower bound on the reference coupling
so each row sum, including the below-diagonal terms, is at most $1/2$.
The finitely many adjacent-stage cases are included before each choice.
This proves a homogeneous contraction of the normalized defect vector.
Starting with the Wilson defect, whose normalized norm is bounded, its
value after $n$ transfers is at most $C2^{-n}$. At a fixed scale
the logarithmic normalization is fixed, so the unnormalized
defect derivatives tend to zero as well.

\subsection{Why action matching introduces no inhomogeneous term}

We specify the map on vector fields, including its sign. At the critical
configuration and coarse local gauge transformation, let $\widetilde Q(U,h)$ be the
averaging operation with its chosen local gauge transformation. Its left-trivialized differential is
\[
 (R_f\eta)_e=\left.\frac d{ds}\right|_0
 \log\bigl(\widetilde Q(e^{s\eta}U,h)_e
                              \widetilde Q(U,h)_e^{-1}\bigr).
\]
Let $b_{y,f}$ be the fixed contour coefficients of the local gauge transformation, defined by
$\sum_f b_{y,f}\eta_f=\tau_y(0;\eta)-\sum_rw_r^{\rm av}\tau_y(r;\eta)$,
where $\tau$ is the signed path average in
\eqref{ch05:frame-linearization}. With the fixed contour $\gamma_{y,f}$
from coarse root $Ly$ to the initial vertex of the fine link $f$, set
\[
 (R_0\eta)(y)=\sum_f b_{y,f}
             \operatorname{Ad}(h(y)U(\gamma_{y,f}))\eta_f,
 \qquad (d_V\phi)_e=\phi(e^-)-\operatorname{Ad}(V_e)\phi(e^+).
\]
At coarse base point zero the transferred coarse vector field is
\begin{equation}
 Y'_b=L^5\sum_xK_0(x)(R_f+d_VR_0)
                              Y_b^x(U_{\rm stat}(V)).
 \label{ch17:actual-current-map}
\end{equation}
Translations give the other base points. The sum is over fine base points;
the output is a vector on coarse links. At the identity $R_0=R$ of
\eqref{ch05:actual-Q}, $R_f=Q=M+d_{0,c}R$, and
$d_V=-d_{0,c}$. In this convention the plus sign in
\eqref{ch17:actual-current-map} cancels the root correction exactly:
\begin{equation}
                         (R_f+d_VR_0)|_{0}=M.
 \label{ch17:flat-current-map}
\end{equation}
The smooth linear averaging operator, rather than a rough row along the root contours,
therefore controls its equal-order Taylor terms. All derivatives of
$R_f+d_VR_0$ have finite $C_{j,L}$ bounds on the common coarse strip.
Their fixed contours have length $O_j(L)$ and remain within bounded
coarse distance; differentiation adds only finite contour-length factors.

For the exact action at the critical configuration the current identity is preserved before
any estimates. The envelope theorem differentiates the stationary value
only through the explicit coarse field. In the penalty term, simultaneous
left multiplication of the averaging operation and the coarse link leaves the
penalty term unchanged. Thus a fine vector field $Y$ induces the coarse
vector field $R_fY$; its root correction $d_VR_0Y$ has zero action
variation by coarse gauge invariance. Together with
\eqref{ch17:block-algebra} this gives
\[
 DS_{\rm stat}[Y'_b]
  =L^5\mathcal I K_0DS[Y_b],\qquad
 \sum_a\nabla_a^{\rm c}J'_{ab}
  =L^5\mathcal I K_0\sum_a\nabla_aJ_{ab}.
\]
Joint stationarity in the fine field and the local gauge transformation is used here, as in the gauge-fixed
quotient calculation of Chapter~\ref{ch07:matching}. Matching replaces
$S_{\rm stat}$ by an action with identical Taylor coefficients through
field order four. Since $Y'(0)=0$, its action variation still agrees
through order four. Therefore the displayed defect pullback is exact
at every field order used in \eqref{ch17:unnormalized-defect}.
A separate additive $C/(\log\widehat\beta)^2$ matching error is neither
needed nor permitted in this argument.

\subsection{Uniform current and difference-primitive bounds}

Recall that $J_{ab}$ is the lattice stress current and $Y_b$ its
associated gauge-covariant vector field. The functions $q_a$ and
$t_{abc}$ are the forward-difference primitives of its trace and
antisymmetric part, respectively, as defined in
Section~\ref{ch17:current-suppliers}. They are distinct from the
small-field threshold $t$ used in the norms with insertions.

For completeness we describe the accompanying bounds needed in the
regulator criterion. Subtract fixed local reference stress and curvature
vector fields, whose leading Taylor coefficients by scaling dimension are those used initially.
The reference functions have finite stencils and fixed coefficient bounds.\mgubi{Issue 8: gauge invariance of the reference stress is not stated; the division by two curls below uses the Ward identity on each leg.}
The lowest jets of the transferred functions agree with them exactly.
For the stress this follows from the exact rescaling at the end of the proof of preservation of the order of vanishing,
applied with $n=2$ to the Hessian $B$ of $K_aJ_{ab}$. Its values on pairs of polynomial fields with $r+r'<2$ vanish: one leg is then constant,
a constant field is the gradient of a linear potential, and the truncation argument given there, using the exponential decay of the kernel, extends the Ward identity to such potentials. If
$r+r'=2$, the rescaling gives $D^2\mathcal I(K_aJ_{ab})(0)[P,P']=L^{-4}B(P,P')$, which the factor $L^4$ in $J'_{ab}$ compensates.
Moreover $K_a$ is a finite combination of fine translates with coefficient sum one, and a translate of a homogeneous polynomial
differs from it by terms of lower degree, on which the Hessian of $J_{ab}$ vanishes; so $B(P,P')$ is the corresponding value for $J_{ab}$.
Consequently the quadratic coefficients of scaling dimension below four remain zero, and the quadratic
coefficient of scaling dimension four is unchanged.

For the vector field one additionally uses the column sum of the linearized averaging operation.
By \eqref{ch17:flat-current-map}, the corrected differential is exactly
$M$ at zero. The alias formula
for $M$ annihilates every nonzero alias at zero coarse momentum, so its
fine column sum in a fixed direction is constant. Its row sum is $L$;
there are $L^4$ positions within a block. The column sum is therefore $L^{-3}$.
This is an identity for the corrected differential, with no residual
root-gradient term. Since the leading fine curvature of $W_a$ is $L^{-2}$ times
the coarse curvature, the factors in the transfer of the current multiply to
$L^5L^{-3}L^{-2}=1$. This preserves the linear curvature jet of $Y$.
Gauge covariance then fixes its accompanying quadratic commutator term:
the covariant polynomial of scaling dimension two in the transferred
vector field is again the curvature.

The stress remainder's quadratic symbol now has order at least three.
Two-leg curl division leaves one extra difference, so its diagonal
transfer is $C_jL^4L^{-5}\le C_j/L$. Higher stress orders have
radial-gauge bounds $C_jL^{4-2r}$.
For $Y$, its zero-order corrected differential $M$ has row sum at most
$C_jL$, while $\|K_0\|_\infty\le C_jL^{-4}$.
To see the resulting local gain, write a translated current coefficient
as $y(f-x;z_1-x,\ldots,z_r-x)$, with output link $f$ and base point $x$.
For each fixed coarse output edge $e$ its absolute sum is bounded by
\[
 \|K_0\|_\infty\sum_f|M_{ef}|
        \sum_x|y(f-x;\,\cdot\,)|
                   \le C_jL^{-3}\|y\|.
\]
Only a bounded number of coarse output edges are within the patches of the averaging
operation and of the base-point average; exponential coefficient tails control the rest.
This proves the same gain for the full local output-link norm.
The linear vector-field remainder has a curvature difference and one
additional difference. If the extra difference acts on its field leg,
it contributes a factor $L^{-1}$ through $W_a$. If it acts on the fine output index,
summation by parts moves it onto $M_{ef}$; the identity for the linear averaging operator gives
$\sum_f|\nabla_fM_{ef}|\le C_j$, again gaining $L^{-1}$ relative to
its row sum $C_jL$. Thus its diagonal factor is
$C_jL^5L^{-3}L^{-3}=C_j/L$.
For field order $r\ge2$ its diagonal factor is $C_jL^{2-2r}$.
Derivatives of the averaging operation or of the correction from its local gauge transformation multiply lower
field orders of $Y$, so their fixed-$L$ constants are absorbed by the
same logarithmic normalization as above. Evaluating in a local radial
gauge is legitimate for this covariant vector field: transform the local gauge
transformation at its output and the differential of the averaging operation together. The additional
derivatives of the local gauge transformation again occur only below the field-order diagonal.

The trace and antisymmetry primitives start at field order two. Their
old terms carry $L^3$ and hence gain $C_jL^{3-2r}\le C_j/L$.
Their new terms are the explicit $L^4K_a^{\rm t}J$ terms proved in
Lemma~\ref{ch17:block-identities}; the equal-order inhomogeneous term from $J$
is bounded independently of $L$. Choose a fixed smaller weight for
these primitive norms so this one-way inhomogeneous term has norm at most $1/8$.
The other diagonal terms are then made small by $L$, and all lower-order
terms by the lower bound on the reference coupling. Reference functions contribute finite
$C_{j,L}$ inhomogeneous terms, divided by positive powers of $\log\widehat\beta$;
a sufficiently large lower bound on the reference coupling can make these terms at most $1/4$.

Thus a weighted sum of the normalized current and primitive norms obeys
$X_m\le\frac12X_{m+1}+\frac14$. The initial Wilson current and zero
primitives have bounded normalized norms. Uniform bounds follow throughout
the history, and in particular at each fixed scale. At that scale
these are the unnormalized finite constants required by the regulator
criterion, since $\widehat\beta_m$ is fixed. The needed free scaling-dimension
jets are unchanged, and the defect tends to zero by the homogeneous
contraction above. All estimates use the same averaging operation, critical
configuration, action matching and periodized kernels as the trajectory with insertions.


\section{Quantitative stability of the trace functional}
\label{ch17:stability}

The finite-dimensional argument is elementary, but its rate is essential.
Suppose $s:E\to V$ has right inverse $t$ and one-dimensional kernel
$\RR v$. Let $s_i:E\to V$, $t_i:V\to E$ satisfy
$s_it_i=I$, $t_i\to t$, and let $v_i\to v$. If $e_i\ge1$ and
\[
 e_i\|s_iv_i\|\to0,\qquad e_i|\rho(v_i)|\to0,
\]
then, for $\lambda_i=\rho t_i$,
\begin{equation}
 s_i\to s,\quad \sup_i\|\lambda_i\|<\infty,
              \quad e_i\|\lambda_is_i-\rho\|\to0.
 \label{ch17:stable-functional}
\end{equation}
Indeed $k_i=v_i-t_is_iv_i\to v$ lies in $\ker s_i$. The maps
$(h,c)\mapsto t_ih+ck_i$ converge to the invertible map
$(h,c)\mapsto th+cv$, so their inverses are uniformly bounded.
Writing $y=t_is_iy+c_i(y)k_i$ gives $|c_i(y)|\le C\|y\|$ and
$(\lambda_is_i-\rho)y=-c_i(y)\rho(k_i)$, which proves the claim.

Here are the complete rate estimates needed for the trace conclusion;
the next section proves them for the common family.
Write
\[
 e_i=\exp((\log\log b_i)^2/5),\qquad
 q_i=\exp((\log\log b_i)^2/6),
\]
so $q_i/e_i\to0$ and both are smaller than every positive power of
$b_i$. The letter $q_i$ in this section is a scalar bound, unrelated
to the inserted densities of Chapter~\ref{ch11:insertions}.
Let $B_i^L,B_i^X,B_i^D$ be the centered transferred tensors and set
\[
 A_i=\beta_i r_i^4 U_{i,4}^{-1},\qquad
 K_i=\beta_i r_i^5 U_{i,5}^{-1}\operatorname{div}.
\]
The required same-family current transfer estimates are
\begin{align}
 &\|\beta_iK_iB_i^L\|\le Cq_i,\quad
             \lambda_i(\beta_iK_iB_i^L)\to-\kappa_G,
 \label{ch17:transfer}\\
 &K_iB_i^Z\to\operatorname{div}^{\rm Max}Z,\quad
 e_i[\lambda_i(K_iB_i^Z)-\rho(\operatorname{div}_{\rm raw}Z)]\to0
                       \quad(Z=X,D),\notag\\
 &A_iB_i^Z\to Z\quad(Z=L,X,D),\quad
 b_i^p\operatorname{tr}(A_iB_i^Z)\to0
                       \quad(Z=L,X,\ p\in\NN).\notag
\end{align}
These bounds are stronger than smooth free comparison. The next section
proves them using the kernel identities and a separate low-moment
estimate for spatial remainders.

% !TEX root = ../monograph.tex
\section{From the current identity to coefficient estimates}
\label{ch17:marked-rates}

The integration by parts needed here is performed before the
large-field expansion. This order matters: a large-field activity is
only measurable, so its derivative cannot be inserted into the
expansion as an ordinary bounded function.

\subsection{A differentiated block kernel}

In the exact kernel $k^{(L)}_\beta$ of \eqref{ch04:kernel}, condition
on the fine field $U$ and all auxiliary choices $\alpha$. The conditional density of $V_e$ is
$q_{b\beta}(V_e Q_e^{-1})$, where
$q_\gamma(u)=Z_\gamma^{-1}\exp[-\gamma w(u)]$.
For every fixed order $r$, compactness and smoothness of $w$ give
\begin{equation}\label{ch17:kernel-score}
 |D^r q_\gamma(u)|\le C_r(1+\gamma)^r q_\gamma(u).
\end{equation}
Here $D^r$ denotes $r$ successive left derivatives in unit Lie algebra
directions. Dividing by $q_\gamma$, the result is a polynomial in
derivatives of $w$ with at most $r$ factors of $\gamma$. The normalization is
independent of $u$. For a product of link kernels the corresponding
bound has an additional polynomial in the number of differentiated
links. No lower bound on $q_\gamma$ is used.

First truncate the exponentially local current and the kernels of the critical-configuration map (the coefficients of its Taylor expansion in the coarse field) on a finite neighborhood. For the vector field use a finite
gauge-covariant extension of its controlled jets through field order four,
not the uncontrolled higher derivatives of a deep transferred current.
Use the coefficientwise gauge-covariant extension: replace each
fixed-order curvature or covariant-difference monomial by its plaquette
and contour-variable representative, and correct its lower jets
triangularly through order four. Each fixed number of link derivatives
adds only a fixed power of the contour length. The extension therefore
has a bounded weighted norm at a strictly narrower exponential width;
it is not an arbitrary jet section on a space of growing dimension.
Complete in that weighted norm first and then truncate. All additional
fixed-order derivative bounds depend only on the controlled jets and
the fixed stage. Extend this local small-field formula smoothly to
the compact product group. One can use a smooth cutoff
supported strictly inside the logarithmic chart, equal to one on the
smaller chart used for Taylor extraction, and then average over the
compact local gauge group. Equivalently extend its finite covariant
holonomy jets. This preserves the required jets and covariance.
The finite derivative factors are inverse powers of the chart width and
polynomials in the support size. They are allowed in the bounds below.
Only after this extension do we integrate by parts. For an exponentially
local current do this first for finite support truncations. Fixed powers
of the support length are absorbed by a strict margin in the
exponential coefficient width. The derivative bounds and tail errors
are consequently summable, so the finite-period identity passes to
the weighted tail limit.

Let $k(U,V)$ be that exact kernel, let $f(U)$ be the entire incoming
density, and put $F(V)=\int k(U,V)f(U)\,dU$. All original fine large-field indicators
and compatibility conditions are part of $f$ or of the integral over the local gauge transformations;
none is differentiated. For a smooth coarse vector field $Y$, define
\[
 \mathcal D_Y k=(\operatorname{div}_{\rm Haar}Y)k+D_Vk[Y].
\]
Dominated differentiation, using \eqref{ch17:kernel-score}, gives the
\emph{density identity}
\begin{equation}\label{ch17:kernel-divergence}
 \operatorname{div}_{\rm Haar}(FY)
       =\int\mathcal D_Y k(U,V)f(U)\,dU.
\end{equation}
Consequently its integral against every spectator independent of the
varied coarse links is zero. This is an identity for the full differentiated kernel,
including its nonpolynomial remainder, not an assertion that a finite
Taylor polynomial alone has zero exterior moments, that is, zero integral against every spectator independent of the varied links.

The right hand side is an admissible conditional integral with the differentiated kernel inserted.
Its extra factor is a finite sum of smooth local kernel scores times
$Y$ or its derivative. On a protected dependence set the incoming
measurable product remains fixed under the change of variables in
Chapter~\ref{ch06:large-fields}. Additional coarse derivatives act
on the score and the common density ratio only. The bounded pointwise-jet estimate
\eqref{ch06:actual-jet-integral} therefore applies with a fixed
polynomial factor in $\beta$, support size and current derivative bounds.
Rooting the connected expansion at this extra factor gives the same
expansion with insertions as in Chapter~\ref{ch11:insertions}. In particular
large-field terms retain their original prescribed large-field sets.
This argument does not differentiate the event indicators introduced when that
conditional integral is subsequently decomposed.

The small-field contribution needs a sharper identity than the absolute
score bound. Fix the full dependence set of a rooted connected component and use one
additional coarse derivative beyond those prescribed for the extracted
coefficient. On its smaller small-field chart, make exactly the conditional
Gaussian change of variables used to construct the small-field expansion.
Write its normalized conditional integrand as $p(V,z)$ and its integral
as $Z_{\rm ord}(V)$. For finite small-field selections, differentiation
under this integral gives
\[
 D_VZ_{\rm ord}[Y]=\int D_Vp[Y]\,dz,
 \qquad
 \int D_Vp[Y]\,dz
 =Z_{\rm ord}D_V\log Z_{\rm ord}[Y].
\]
Every Jacobian, relative determinant and conditional normalization is
included in $p$; thus this is not the derivative of an unnormalized
Gaussian moment. The absolutely convergent small-field logarithm, with
one additional derivative, passes this identity to the countable expansion.
By its definition in Chapter~\ref{ch09:renormalization},\mgubi{Issue 9: $Z_{\rm ord}$ is defined only above, for one component; $E_{\rm reg}$ is defined nowhere. Chapter~\ref{ch09:renormalization} gives this form with a constant only for related objects (the pure small-field block integral; the constant extracted once, globally; part (a) of Assumption~\ref{ch09:forest-assumption}, unproved). A constant for one component needs an argument.}
\[
 Z_{\rm ord}=\mathcal C
       \exp[-\beta S-\mathcal L-E_{\rm reg}],
 \qquad D_V\log\mathcal C=0.
\]
If an intermediate local normalization depends on $V$, it is included
in $\mathcal L+E_{\rm reg}$, not in $\mathcal C$. Consequently the
small-field score is \emph{exactly}
$-\beta DS[Y]-D\mathcal L[Y]-DE_{\rm reg}[Y]$.
The regular term is estimated using its already proved differentiated
regular norm. We do not estimate it by the polynomial-in-$\beta$
absolute score bound.

This argument uses the full rooted component, not just the score
stencil. In the rooted connected expansion, small-field factors in
components disjoint from the component of the differentiated kernel cancel
between the differentiated numerator and its normalization. Every term with a large fluctuation-field event anywhere in the connected
component of the differentiated kernel, including a distant such event joined to the score, is kept in the complementary
sum. The smooth cutoffs used above are constant on the smaller small-field
chart; their derivatives are supported where at least one enlarged
chart condition fails. Attach this condition to its dependence set. The same
Gaussian or energy estimate for the nested cutoffs bounds this new
large fluctuation-field event, while all incoming prescribed large-field sets are retained.
The absolute score bound is used only for these complementary terms,
where the exponentially small probability of the large fluctuation-field event absorbs its polynomial
$\beta$ factor. No derivative of a hard indicator is taken. This proves
the small-field cancellation and the complementary large-field estimate
separately, which is essential to preserve the regular power gain.

\subsection{The finite degree-five coefficient}

At the comparison scale $k$ write
\[
 h=(\log\log\widehat\beta_k)^2,\qquad H=e^h,
 \qquad t=\widehat\beta_k^{-1/2}H.
\]
All constants in this paragraph are for a fixed stage and the fixed
finite derivative orders and degrees; they are fixed before its lower bound on the
reference coupling. Let $\mathcal P_5$ denote radial Taylor extraction through
degree five and let $f$ be the number of curvature factors of a homogeneous
coefficient. Set $\mathcal P_{5,t}=t^{f-2}\mathcal P_5$, followed by
real signed-permutation averaging. There are no invariant linear forms on $\gLie$. Thus the nonzero vector coefficients under consideration
have $f\ge2$, and this normalization does not introduce a negative
power of $t$ beyond the displayed $t^{-2}$.

Apply the exact step with insertions to \eqref{ch17:kernel-divergence}, retaining
its entire remainder, and use the current identity through order four.
Let $j_k$ be the radial coefficient of the current divergence, in the
same threshold coordinates. The classical term $-\beta DS[Y]$ is
$-\beta j_k$ up to the contracting current defect. Thus the coefficient
of the full differentiated kernel, with zero exterior moments, is $\Sigma_k-\beta j_k$;
$\Sigma_k$ alone is not asserted to have zero exterior moments. Its
remaining degree-five scalar-vector coefficient can be written
\begin{equation}\label{ch17:finite-sigma}
 \Sigma_k=\mathcal P_{5,t}
       (\operatorname{div}_{\rm Haar}Y-D\mathcal L[Y])
       +E_{\rm stat}+E_{\rm tail}+E_{\rm reg}+E_{\rm mark}.
\end{equation}
Here the kernels of the critical-configuration map and the loop kernels are truncated at
$C_j(\log\widehat\beta_k)^2$. Their common exponential decay makes
$E_{\rm stat}$ and $E_{\rm tail}$ at most $H^{-1}$ after finite
projection factors. The regular error satisfies
\[
 t^{-2}\epsilon_k^{\rm reg}
       =\widehat\beta_k^{-1/10}H^{-2}.
\]
A directional derivative along $Y$ has the same bound up to the fixed
current norm, since the derivative bound contributes a factor $t^{-1}$
and $Y=O(t)$ on the small-field chart. The previously proved logarithmic current bounds
and fixed localization factors are absorbed by increasing the lower bound on the reference coupling.

For the remainder with insertions one must first cancel the whole classical
current divergence, not merely use the absolute kernel-score bound.
The current identity agrees through field order four. The finite
gauge-covariant extension just chosen has controlled higher derivatives,
so the remaining classical error starts at order five and contributes
\[
 t^{-2}\beta t^5=O(\widehat\beta_k^{-1/2}H^3).
\]
Fixed current and localization factors preserve a bound of the form
$C_j\widehat\beta_k^{-1/10}H^{C_j}$ after increasing the lower bound on the reference coupling. The small-field loop
and Jacobian Taylor remainders have no leading factor $\beta$ and
gain an additional small-field factor. This controls the high-field-order remainder. It does not control a
low-field-order term with many spatial differences: vanishing Taylor coefficients through the specified scaling dimension alone give no additional small-field factor. The remaining low-field-order spatial terms are controlled by exact
polynomial preservation, rather than by an extra power of the field.

\begin{proposition}[Low-moment estimate for spatial remainders]
\label{ch17:spatial-moments}
For the linearization of the critical configuration and the regular local terms of the differentiated kernel just constructed,
a remainder whose radial Taylor coefficients through scaling dimension five vanish
has extracted degree-five coefficient $O(H^{-1/2})$. The estimate is
uniform in period and in the finite current truncations, after the
stage orders and thresholds have been chosen with one extra derivative
order. It applies to both the nonclassical coefficient and the
normalized current divergence.
\end{proposition}
\begin{proof}
First identify the linearization of the critical configuration on polynomials. Recall the notation of
Chapter~\ref{ch05:gaussian}: $D_l,D$ are the rescaled fine and coarse
covariance symbols, $x_l,x$ their gradient columns, and
$\alpha=\sum_l|m_l|^2\tr D_l$.
Write $s=\lambda_{M_c}(p)$ and
$D=s^{-1}\operatorname{diag}(B_a)$, where the $B_a$ are the
regularized diagonal covariance factors. The scalar correction
\[
 \chi=\frac{x^*\operatorname{diag}(B_a^{-1})x}{s}
\]
is analytic and nonzero at zero, as in the pole cancellation for the
formulas for the linearization. Put
$\sigma=\chi^{-1}x^*\operatorname{diag}(B_a^{-1})$, so that
$\pi=s^{-1}\sigma$ in those formulas. For a nonzero alias define
\[
 g_{l,a}=\overline{m_l}
 \left(\frac{\tr D_l}{s\alpha}
       -\frac{(D_l)_{aa}|r_L(\theta_{l,a})|^2}{B_a}\right).
\]
Substitution into the linearization built with the linear averaging operator gives the regular formula
\[
 \mathcal A_l=\mathcal B_l+
          \operatorname{diag}(g_{l,a})x_l\sigma,
 \qquad \mathcal B_l=D_lF_{1,l}^*D^{-1}.
\]
Here $D_l$ is regular at a nonzero alias, and $s\alpha$ and
$B_a$ are analytic and nonzero at zero. Thus no uncancelled factor
$s^{-1}$ remains. Both terms contain $\overline{m_l}$, with
$m_l=w_L(\theta_l)^j$ and $w_L(\theta_l(0))=0$.
Every derivative of $\mathcal A_l$ of order less than $j$ therefore
vanishes at zero. The column at the zero alias has value $I$. Through \eqref{ch17:polynomial-response}, this
order-$j$ vanishing is also what the preservation of the order of vanishing of the
current defect in Section~\ref{ch17:quadratic-current} uses. Stars and conjugates in these symbols mean analytic
transpose at the opposite momentum, as in Chapter~\ref{ch05:gaussian}.

Let $A_x^{ab}(z)$ be the physical row from coarse direction $b$ at
$\lfloor x/L\rfloor+z$ to fine direction $a$ at $x$.
Differentiating the Fourier inversion formula through order less than $j$
therefore gives, for every vector polynomial $p$ of degree less than $j$,
\begin{equation}\label{ch17:polynomial-response}
 (W_sp)(x)_a=\sum_{b,z}A_0^{ab}(z)p_b(x/L+z),
 \qquad \sum_z A_0^{ab}(z)=L^{-1}\delta_{ab}.
\end{equation}
The exponential row bound justifies the derivatives and the sum.
Polarization gives mixed polynomial moments. In particular the output
is a global polynomial of no larger degree, not a polynomial separately
on each block. Its homogeneous degree-$r$ leading term is
$L^{-r-1}p(x)$; lower coefficients obey the same ratio-independent row bounds.

Choose $j$ above the order of the coefficient evaluations. In a polarized coefficient
evaluation take only tuples of potential polynomials whose summed scaling
dimension is at most five. Equation~\eqref{ch17:polynomial-response}
preserves or lowers each polynomial degree, so every resulting tuple
remains within that same finite list. The polynomial-moment
cancellation of Section~\ref{ch11:one-graph} makes its deterministic
coefficient evaluation exactly zero. This is not a claim of vanishing on arbitrary
polynomial inputs of degree less than $j$. Truncating the coarse field to the ball of radius
$C_j(\log\widehat\beta_k)^2$ changes this coefficient evaluation only by the
exponential row tail, times fixed polynomial directional and support
factors. That tail is smaller than every prescribed power of $H^{-1}$.

The smooth representative of the linearization of the critical configuration is
$W_a=W_s-d_fERW_s$. Its difference from $W_s$ is a fine gradient.
The coefficient of these regular terms is a spacetime vector with trivial adjoint action on its values,
invariant under local fine gauge transformations. Indeed gauge
equivariance of $Y$ and invariance of Haar volume imply invariance of
$\operatorname{div}_{\rm Haar}Y$; the action and loop variations
are invariant for the same reason. Apply the corresponding fine gauge
transformation before evaluating its local jet. The Baker--Campbell--Hausdorff formula with the same local gauge choice compares the two exponential
representatives with an error $C_jP(B)t^2$ in the link logarithms,
where $P$ is a fixed polynomial. Dividing by the fine threshold gives
$C_jP(B)t$. This argument uses the exact correction to the averaging operation, not an
unmatched linear averaging operator.

It remains to compare the conditional insertion itself to this
polynomial coefficient evaluation. Normalize the input link logarithms by $t$.
The nonlinear critical configuration, the same local gauge coordinate change,
and the deterministic truncation contribute $C_jP(B)t$ and the exponential
tail just described. Gaussian fluctuations contribute
$C_jP(B)g/t$, where $g=\beta^{-1/2}$, by the fundamental theorem of
calculus and the one additional derivative. Explicitly, if
$p(z)\ge1$ is a polynomial majorant for the finite Gaussian coordinates,
the difference between the conditional and the deterministic evaluation in
each prescribed derivative is bounded by
\[
 C_j N_k P(B)\bigl(t+g/t+e^{-c_jB}\bigr)\,p(z)^{C_j},
\]
with an appropriate enlargement of the truncation radius in the last term.
Here $N_k$ is the norm of the already $t^{-2}$-normalized regular
term of the differentiated kernel and its derivatives in normalized link logarithms, including
its current and loop coefficients. The extraction normalized by the number of curvature factors has
already been performed; no further $t^{-2}$ is applied to this bound. For this bound subtract the deterministic coefficient evaluation
inside the same normalized conditional integral. Its derivative
through the prescribed scaling dimension is zero. Use the joint
polynomial moments and differentiated density scores of that one
conditional law to estimate the evaluation difference; no small
total-variation distance between the nonlinear law and a Gaussian
is asserted. The pointwise jet and moment estimates are those of \eqref{ch06:actual-jet-integral};
no differentiation of incoming large-field factors is involved.
The complement is estimated by the already retained large-field bound. Finite coefficient recovery adds a factor depending only on the stage.

There are no invariant terms that are constant spacetime vectors or linear in the $\gLie$-valued field. Hence these regular terms of the differentiated kernel start at field order two. Their
values and normalized-link derivatives on the threshold ball contain
a factor $t^2$; the explicit $t^{-2}$ removes exactly this factor.
The classical part has already been cancelled and its order-five
remainder treated separately above. Thus for the present regular terms
$N_k\le H^{\theta_j+o(1)}$, with
$\theta_j\to0$: current norms are logarithmic and the retained loop
norm has this bound. Since $B$ and all truncation radii are fixed powers
of $\log\widehat\beta_k$, their polynomial factors are $H^{o(1)}$.
The coupling band gives $g/t\le C H^{-1}$, whereas $t$ and
the truncation tail decay faster than every fixed inverse power of $H$.
The resulting coefficient error is $H^{-1+\theta_j+o(1)}$, hence
at most $CH^{-1/2}$ after increasing the lower bound on the reference coupling. Large supports are first truncated
using a margin in the support exponent; its tail obeys the same
bound. This proves the proposition for the full local sum.
\end{proof}

The differentiated-kernel
bound \eqref{ch17:kernel-score} applies to the complementary part.
The complementary part retains the square root
$\exp(-H^{1/5}/2)$ of the bound \eqref{ch04:activity-bounds} on the large-field activities. Combining Proposition~\ref{ch17:spatial-moments} with the high-field
remainder, the nonwrapping contributions are bounded by
\begin{equation}\label{ch17:sigma-error}
 C H^{-1/2}+C_j\widehat\beta_k^{-1/10}H^{C_j}
       +C_j\widehat\beta_k^{C_j}H^{C_j}e^{-H^{1/5}/2}.
\end{equation}
The fixed-order score has changed the polynomial constants, not the
eligibility event \eqref{ch04:eligibility} of a large-field activity or its exponential gain. If $N_{\rm per}$ is the
period in scale-$k$ units, a wrapping term has covering number at least
$N_{\rm per}/[10(2B_k+1)]-1$. Retaining two units of the support exponent gives the
additional error
\[
 C_j\widehat\beta_k^{C_j}H^{C_j}
       \exp\bigl(-N_{\rm per}/[5(2B_k+1)]+2\bigr).
\]
It vanishes in the ordered period limit at fixed $k$.

The two terms other than $CH^{-1/2}$ in
\eqref{ch17:sigma-error} are $o(H^{-1})$. To check the
order of limits explicitly, put $x=\log\log\widehat\beta_k$.
The logarithm of the first bound multiplied by $H$ is at most
$C_j+(C_j+1)x^2-e^x/10$; that of the second is at most
$C_j+C_je^x+(C_j+1)x^2-e^{x^2/5}/2$.
Both tend to minus infinity at each fixed stage. Increasing its lower bound on the
reference coupling makes the bounds uniform on that stage.

The loop coefficient has zero $\rho$ by its two moment cancellations:
its quadratic kernel annihilates both constant and affine links, whereas
the cubic momentum coefficient functional pairs a constant with a quadratic link or two
affine links. This uses the loop kernel of Chapter~\ref{ch07:matching},
not only gauge invariance.
The calculation with mass regularization identifies the remaining coefficient as
$\mathfrak h(Y)/12$. One can check the factor directly in the three
vector coordinates of \eqref{ch17:rho}, before the quotient by the equation of motion: the cubic momentum
functional takes values $-24,-96,0$, while $\rho$ takes values
$-2,-8,0$. The lattice correction is the matching lower-order
Taylor correction in $\mathfrak h$.

Choose the number of transfers and then the period at each fixed $k$ so
that the error of the mass regularization and the finite momentum approximation are
smaller than $\widehat\beta_k^{-30}$. The retained loop bound is
$2H^{\theta_j}$, with $\theta_j\to0$, and current derivative norms
are logarithmic. Sufficiently large lower bounds on the reference coupling therefore give joint ordered limit
coefficients satisfying
\begin{equation}\label{ch17:sigma-bound}
 \|\Sigma_k\|\le C H^{1/6},\qquad
 \rho(\Sigma_k)\longrightarrow-\kappa_G.
\end{equation}
All coefficient maps and the underlying correlations are retained
in the same joint limit. Compactness is used only after the uniform
finite-cutoff estimates have been established.

\subsection{The kernel vector and the necessary rates}

Let $j_k$ be the degree-five coefficient of the transferred current's
divergence, in the same threshold coordinates. Its local Taylor
coefficient has zero $\rho$: the trace is a forward difference of a
degree-three primitive, whose invariant projection is a multiple of the identity field.
The same localization and error bounds with insertions compare the transferred
coefficient with this Taylor coefficient within $CH^{-1/4}$.
Consequently, in the notation of the stability argument, with $v=(0,0,1)$ the kernel vector of the free quotient $s$,
\[
 j_i\to-v,\quad |\rho(j_i)|\le CH_i^{-1/4},\qquad
 v_i^{\rm eq}=\beta_i^{-1}\Sigma_{k_i}-j_i\to v.
\]
With $e_i=H_i^{1/5}$ this gives $e_i\rho(v_i^{\rm eq})\to0$.

There is also a stronger, independent bound on the image of this
vector. The differentiated kernel is carried as a density-valued factor, without
dividing it by $F$ or claiming that the resulting quotient is local.
More directly, for every incoming coefficient $f_A(U)$ with insertions,
including all the original fine spectator factors,
\[
 \int \mathcal D_Y k(U,V)\,dV=0
 \quad\hbox{for each }U,\qquad
 \int \mathcal D_Y k(U,V)f_A(U)\,dU\,dV=0.
\]
The same differentiated kernel is used for every $A$, so the joint rooted
expansion retains all pushed spectator cross-modes and this exact
zero-moment statement. There is no smoothness requirement on $f_A$.
For a smooth local coarse spectator $A(V)$ one also has
\[
 \int A(V)\!\int\mathcal D_Y k(U,V)f(U)\,dU\,dV
 =-\int D_VA[Y]\!\int k(U,V)f(U)\,dU\,dV.
\]
It is zero when the spectator does not read the varied links. For an
exponentially local spectator that need not have finite range, first truncate its dependence set
and use this identity on disjoint terms. Retain the remaining terms
in the same joint expansion with insertions. A connected term reaching from
the differentiated kernel to a spectator at distance $R$ has covering number at least
$R/(2B_k+1)-1$. Retaining two units of the support exponent bounds these terms by
$C\exp[-2R/(2B_k+1)]$, with their normalized large-field majorants as well.
For spectator configurations separated by a fixed positive multiple of $a_n$, their
distance in scale-$k(n)$ units is at least $c\exp(T_n)$. This tail
is therefore smaller than $e^{-T_n/100}$ after all fixed factors.
Every spectator subset uses the same density without insertions.
Write $\tau_j=\log(1/a_j)$ and
$T_n=(\log\widehat\beta_n)^{48}$. Choose $k(n)$ to be the smallest
integer $k\ge n$ for which $\tau_k-\tau_n\ge T_n$. Match the full
coefficient and remainder backwards from $k(n)$ to $n$. Equation~\eqref{ch11:descent-error} gives
$e^{-T_n/8}$ before multiplication by its constants. Fixed projection and
threshold factors have logarithm $O(\log\widehat\beta_n)$; accumulated
ratio-independent coefficient-map factors are $o(T_n)$. Here it is only the latter
constants that are multiplied along the interval of scales. Ratio-dependent nonlinear
jet constants have already been absorbed into the errors controlled by the lower bounds on the reference coupling of the stages.

Place a copy in each reflected cell at the first backward crossing of $B_n^2$, so
$B_n^2\le a_m/a_n\le L_*B_n^2$. The lower bound on the partition function
in \eqref{ch11:replicated-bound} contributes at most
\[
 L_*^4B_n^8q_n+6=O((\log\widehat\beta_n)^{22})=o(T_n),
\]
using $L_*\le\log\widehat\beta_n$, $B_n=O((\log\widehat\beta_n)^2)$
and the per-site normalization exponent $q_n=O((\log\widehat\beta_n)^2)$.
The inverse source weights in \eqref{ch11:replicated-bound} have logarithm $O(\log\widehat\beta_n)$.
The dimensionful normalization of the insertion cancels the factor $a_n^{-5}$.
Thus the normalized residual of \eqref{ch11:descent-error} is at most $e^{-T_n/100}$.
This argument uses the full differentiated kernel with zero exterior moments; its Taylor remainder
is carried by the graph of the backward comparison, not discarded.

In the choice of radii along the diagonal subsequence require
$r_i/a_{n_i}\le\exp((\log\log\widehat\beta_{n_i})^3)$.
Finite separating matrices of free two-point correlation pairings have uniformly bounded left
inverses. For every fixed $p$ the image therefore satisfies
\[
 b_i^p\|s_i(v_i^{\rm eq})\|
 \le C b_i^p\exp\bigl(5(\log\log\widehat\beta_{n_i})^3-T_{n_i}/100\bigr)
 \longrightarrow0.
\]
This proves the quantitative kernel hypothesis in the stability lemma.
Its bounded right inverses give bounded $\lambda_i$, with
$e_i\|\lambda_i s_i-\rho\|\to0$.

The same identity for the differentiated kernel gives
$\beta_iK_iB_i^L-s_i\Sigma_{k_i}\to0$.
For $Z=X,D$, the ordinary backward comparison gives
$K_iB_i^Z-s_i\operatorname{div}_{\rm raw}Z=o(b_i^{-p})$ for every
fixed $p$. Apply the bounded functional and use
$H_i^{1/6}/e_i\to0$ in \eqref{ch17:sigma-bound}.
This proves the first two lines of \eqref{ch17:transfer}.

Finally the trace of $X$ is identically zero, while the transferred
trace of $L$ is an exact forward difference. Match its degree-three
primitive at a fixed late pair of cell scale and stopping scale. The backward comparison gives endpoint
error $C e^{-(\tau_k-\tau_{k_0})/8}$, uniformly near the configurations used in the finite separating
correlation pairings. Subexponential coefficient conversions and support truncation
make this $O(a_k^\eta)$ for any fixed $0<\eta<1/8$.
Average over $m=\lfloor a_k^{-1+\eta/2}\rfloor$ consecutive differences.
The identity-field parts of the two endpoint values are equal and cancel, leaving at most
$2Ca_k^\eta/(ma_k)=O(a_k^{\eta/2})$.
The averaging length $ma_k$ in continuum units tends to zero, so the finite matrix of
correlation pairings converges to its original invertible matrix. This proves the
last rate in \eqref{ch17:transfer}, including multiplication by every
fixed power of $b_i$. Ordinary Taylor matching gives its remaining
convergence $A_iB_i^Z\to Z$.


\section{The coefficient and the all-radius limit}
\label{ch17:limit}

\begin{proposition}[The trace coefficient]
\label{ch17:coefficient}
For the prescribed-radius comparison maps constructed above,
\[
          \beta_i r_i^4 U_{i,4}^{-1}\Theta
                       \longrightarrow-\kappa_G\mathcal O.
\]
\end{proposition}
\begin{proof}
The free divergences of $X,D$ form a basis of $V$. Thus the maps
$f_i(c)=c_0K_iB_i^X+c_1K_iB_i^D$ have bounded inverses eventually.
Put $c_i=-f_i^{-1}(\beta_iK_iB_i^L)$, so $|c_i|\le Cq_i$ and
$Z_i=\beta_iB_i^L+c_{i,0}B_i^X+c_{i,1}B_i^D$ is conserved.
Applying $\lambda_i$ to its zero divergence gives
\[
 4c_{i,1}=\lambda_i(\beta_iK_iB_i^L)
 +c_{i,0}\lambda_i(K_iB_i^X)
 +c_{i,1}[\lambda_i(K_iB_i^D)+4].
\]
The last two terms are $O(q_i)o(e_i^{-1})=o(1)$, whence
$c_{i,1}\to-\kappa_G/4$. Since $q_i/\beta_i\to0$,
$r_i^4U_{i,4}^{-1}Z_i\to L$. The faster trace bounds give
$\beta_i\operatorname{tr}(r_i^4U_{i,4}^{-1}Z_i)
\to-\kappa_G\mathcal O$.

The conserved line gives $Z_i=\nu_iT$. Choose one free translation
flux pairing with nonzero value. Its flux on the rescaled $Z_i$ tends
to that value. The Ward identity for $T$ and derivative comparison
make its corresponding flux tend to the same nonzero value.
Their ratio is $\nu_i$, so $\nu_i\to1$. Finally centering subtracts
$e_{\rm tr}\beta_i r_i^4 1\to0$ by \eqref{ch17:running-radius}.
\end{proof}

\begin{theorem}[The all-radius trace law]
\label{ch17:all-radii}
The centered trace in the selected phase satisfies
\begin{equation}
 \lim_{r\downarrow0}r^8\log^2(1/r)G_\Theta(r)
              =\mathcal A_G:=\frac{3\dim_\RR\gLie}{\pi^4}>0.
 \label{ch17:trace-law}
\end{equation}
\end{theorem}
\begin{proof}
The free normalization is
$S_2^{\rm Max}(\mathcal O,\mathcal O;x,0)=\mathcal A_G|x|^{-8}$.
One can check its factor without a tensor convention ambiguity:
the curvature covariance is
$(H_{ai}H_{bj}-H_{aj}H_{bi})/(\pi^2|x|^4)$ with
$H=I-2xx^{\mathsf T}/|x|^2$ orthogonal. The sum of its squared
numerators over ordered pairs is $24$. The two Wick pairings and
the two factors $1/4$ give $2\cdot24/16=3$ per colour.

Start with any prescribed null sequence of radii, and take its
comparison subsequence. Proposition~\ref{ch17:coefficient} gives
convergent moving labels, while \eqref{ch17:running-radius} cancels
the factor $\kappa_G^2$ in their pair correlation. For a generic
pair at distance $t$, lattice rounding therefore yields
\[
 r_i^8\log^2(1/r_i)G_\Theta(r_it_i)\to\mathcal A_Gt^{-8},
                         \qquad t_i\to t.
\]
Choose the pair $0,(t/2,t/2,t/2,t/2)$ for $1/2<t<2$.
For $t<1$ the monotonicity of \eqref{ch17:laplace} bounds the upper
limit at $r_i$ by $\mathcal A_Gt^{-8}$; $t>1$ bounds the lower
limit. Send $t$ to one from each side. Every prescribed null
sequence has a subsequence with limit $\mathcal A_G$, which proves
\eqref{ch17:trace-law} by the subsequence criterion.
\end{proof}

\section{What the trace law excludes}
\label{ch17:exclusion}

Specify the comparison class before using this asymptotic. Consider finite
linear combinations of finite products of the mass-shell Laplace kernels
\[
 K_{m,k,h}(r)=\int_{\RR^3}
 \frac{\omega_m(\boldsymbol p)^k h(\boldsymbol p)}{2\omega_m(\boldsymbol p)}
 e^{-r\omega_m(\boldsymbol p)}\,d^3\boldsymbol p,
 \qquad\omega_m=(m^2+|\boldsymbol p|^2)^{1/2},
\]
where $m\ge0$, $k$ is a nonnegative integer and $h$ is a polynomial;
empty products are constants. The measure has no implicit Fourier
normalization. All these integrals are locally integrable at zero
and converge for $r>0$.

Each such profile has, near zero, a finite sum of terms
$r^a(\log r)^b$, with integer $a$ and nonnegative integer $b$,
plus a bounded remainder. This is the standard Laplace asymptotic expansion
\cite{Wong2001}, applied as follows: decompose $h$ into homogeneous terms,
use polar coordinates, and expand
$(\omega^2-m^2)^{(j+1)/2}$ for $\omega\ge2m$.
Choose enough terms that the remainder has as many energy moments
as needed for the other factors in a product. The retained integer
powers integrate to powers and logarithms; the bounded energy
interval contributes an analytic function. Thus only nonnegative
powers of the logarithm occur.

After multiplication by $r^8\log^2(1/r)$, a nonzero leading term
with $a\le-8$ diverges in magnitude, while every term with $a>-8$
and the bounded remainder tend to zero. Therefore a finite limit,
if it exists, is zero. The positive limit in
\eqref{ch17:trace-law} proves that $G_\Theta$ cannot agree with any
such profile on any interval $(0,\delta)$.

This excludes that specified finite mass-shell polynomial comparison.
It does not exclude a generalized free field chosen to have the
same positive measure $\rho_\Theta$, nor an infinite Wick expansion.
No mass gap follows from a short-distance trace law.
